\chapter{泄流、辐射与热力学关系}

\section{小孔泄流中的能量选择}
\label{th:effusion}
第一讲求出的逸出速率密度为 $g_e(v)=2\alpha^2v^3\ee^{-\alpha v^2}$，$\alpha=m/(2\kb T)$。多出的速率因子来自单位时间的穿孔机会。利用 $I_5/I_3$，有
\[
 \langle v^2\rangle_e=\frac{I_5}{I_3}=\frac2\alpha,
 \qquad \langle E_{\rm tr}\rangle_e=2\kb T.
\]
容器中的平均平动能为 $3\kb T/2$，逸出粒子平均多携带 $\kb T/2$。如果分子还有充分激发的转动自由度，而孔不按内部状态筛选，那么每个内部二次项携带的平均能量仍为 $\kb T/2$。设总有效自由度为 $f$，平动占3个，则
\begin{equation}\label{eq:escapeenergy}
 c_v=\frac f2\kb,\qquad
 \varepsilon_e=2\kb T+\frac{f-3}{2}\kb T
 =\frac{f+1}{2}\kb T.
\end{equation}
这里假定内部自由度与平动保持热平衡，且孔的传输不依赖内部能级。常数 $f$ 只计入所讨论温区内已充分激发的自由度；振动逐渐激发的温区应使用随温度变化的热容。

泄流冷却还需要三个尺度条件：孔径远小于平均自由程，孔外压强足够低，容器内重新平衡的时间远短于粒子数显著变化的时间。第一项保证分子流穿孔，第三项使余气仍可用一个温度描述。它们可以同时满足：小孔很小，容器内部仍有足够碰撞。

\subsection{系统边界与粒子携带的能量}
\label{th:openboundary}
对开口容器应用第一定律时，取固定区域内的瞬时存量作为研究对象。设容器中有 $N$ 个粒子，其内部状态可由 $N,T$ 描述，内能为 $U(N,T)$，短时间内逸出粒子数为 $\dd N_{\rm out}>0$，每粒子平均携带能量为 $\varepsilon_{\rm out}$，则容器粒子数变化为 $\dd N=-\dd N_{\rm out}$。若外界还供热 $\delta Q$，边界无宏观机械位移，
\[
 \dd U=\delta Q-\varepsilon_{\rm out}\dd N_{\rm out}
 =\delta Q+\varepsilon_{\rm out}\dd N.
\]
因 $\dd N<0$，后一项为负。对理想独立粒子，$U=Nu(T)$，完整微分为
\[
 \dd U=Nc_v(T)\dd T+u(T)\dd N.
\]
其中 $u\,\dd N$ 表示存量减少本身造成的内能变化，$Nc_v\dd T$ 表示余气温度改变。与边界输运式联立，
\[
 Nc_v(T)\dd T
 =\delta Q+\bigl[\varepsilon_{\rm out}-u(T)\bigr]\dd N.
\]
余气升温或冷却由逸出能量相对于体内平均能量的高低决定。能量零点若平移 $E_\ast$，$u$ 和 $\varepsilon_{\rm out}$ 同时增加 $E_\ast$，差值不变，温度方程也不变。

\Diagram{fig:controlboundary}{固定边界内的存量与穿越边界的能量}{
 \draw[curve](0,0)--(0,2)--(3.4,2)--(3.4,1.15);
 \draw[curve](3.4,.85)--(3.4,0)--(0,0);
 \node at(1.7,1){$U(N,T)$};
 \draw[axis](-1.6,1)--(0,1)node[midway,above]{$\delta Q$};
 \draw[axis](3.4,1)--(5.6,1)node[right]{$\varepsilon_{\rm out}\dd N_{\rm out}$};
 \node at(2,-.55){先选边界，再写储能与输运};
}
图\ref{fig:controlboundary}中，边界能流按穿越的粒子统计。小孔分子流按穿孔事件取样，平动部分为 $2\kb T$，内部自由度若与穿孔概率无关，则仍取体内平均值。这一微观输运模型与有整体流速的连续稳流模型所用能量形式不同。

对一段连续流管，把入口前的流体推入控制体，外界除送入该流体本身的内能外，还需克服入口压强做功。质量 $\dd m$ 占据体积 $\dd m/\rho$，推入功为 $p\,\dd m/\rho$。因此入口携带的内能加流动功合成比焓 $h=u_m+p/\rho$。在稳定流动中，
\[
 \dot Q-\dot W_{\rm shaft}
 =\dot m\left[
 h_2-h_1+\frac{V_2^2-V_1^2}{2}
 +g(z_2-z_1)\right].
\]
这里 $u_m,h$ 均为单位质量的量，$V_i$ 为宏观流速；$\dot W_{\rm shaft}$ 是对外输出的轴功。入口、出口流动功已经计入焓，不能再作为额外项重复加入。分子流小孔则从单粒子动能与内部能量直接计算输运率；若已这样计入穿面能量，再机械地附加一个 $pV$ 项，便改变了模型。

\subsection{流出冷却与稳流节流的不同约束}
刚性容器缓慢泄流时，容器存量随时间改变，能量方程中有 $\dd U/\dd t$。绝热节流则通常选取阀门附近的固定控制体，在稳态下控制体储能不变；无轴功且宏观动能、位能变化可忽略时，得到 $h_1=h_2$。自由膨胀又是封闭气体占据更大可达体积，外边界刚性绝热，给出整个气体的 $U_1=U_2$。三者分别固定存量能量的变化式、入口出口焓差和整体内能差。

这些约束须在状态方程之前确定。理想气体自由膨胀或理想稳流节流的初末温度相同，是因为它的内能或焓只依赖温度；绝热泄流容器却因穿孔样本偏向高能粒子而冷却。对实际气体，等内能与等焓曲线一般不同，两种温变也须分别计算。

\subsection{持续充气中的流动功}
\begin{example}[绝热刚性空罐的充气温度]
\label{ex:tankfilling}
一个初始接近真空、体积为 $V$ 的刚性容器，通过阀门与温度 $T_s$、压强 $p_s$ 恒定的大型气源相连。气体为常热容理想气体，摩尔热容为 $c_{V,m}$、$c_{p,m}$，$\gamma=c_{p,m}/c_{V,m}$。容器与外界绝热，忽略罐壁热容和最终宏观动能。充气后关阀，待罐内达到均匀平衡。求终温，并说明气源做的功进入了哪里。
\Diagram{fig:tankfillingq}{大型气源经阀门向绝热刚性罐充气}{
 \node[draw,minimum width=2.4cm,minimum height=1.6cm](s)at(0,1){$p_s,T_s$};
 \draw[curve](3,0)rectangle(6.4,2);
 \node at(4.7,1){$V,\ \nu(t),T(t)$};
 \draw[axis](1.2,1)--(3,1);
 \draw(2,.75)--(2.4,1.25);\draw(2,1.25)--(2.4,.75);
 \node[below]at(2.2,.5){阀};
 \node at(4.7,-.5){外壁不动，$Q=0$};
}
\end{example}
\begin{solution}
在图\ref{fig:tankfillingq}中，把连接段和容器共同纳入控制体，在气源侧选入口。每进入 $\dd\nu$ 摩尔气体，气源送入摩尔内能 $c_{V,m}T_s$，并通过入口压强做流动功 $RT_s$，合计送入摩尔焓 $c_{p,m}T_s$。若所选入口处流速不可忽略，应计其动能；本题可在大型气源内部流速可忽略处取截面，直接使用源的摩尔焓。
\[
 \dd U=c_{p,m}T_s\,\dd\nu.
\]
从初始空罐积分，$U_f-U_i=c_{p,m}T_s\int_0^{\nu_f}\dd\nu=c_{p,m}T_s[\nu]_0^{\nu_f}=\nu_fc_{p,m}T_s$，其中$U_i=0$。气体最终静止，内能为 $\nu_fc_{V,m}T_f$，于是
\[
 T_f=\gamma T_s.
\]
\Diagram{fig:tankfillinga}{进入罐内的焓分解为原有内能与入口流动功}{
 \node[draw](a)at(0,0){$c_{p,m}T_s$};
 \node[draw](b)at(4,1){$c_{V,m}T_s$};
 \node[draw](c)at(4,-1){$RT_s$};
 \draw[axis](a)--(b);\draw[axis](a)--(c);
 \node[right]at(6,0){均进入最终内能};
}
图\ref{fig:tankfillinga}中，流动功经过射流、黏滞耗散和内部混合，最终转成罐内热运动能。外壁不动只排除了外壁体积功，入口处的压强仍在输送能量。若充气持续到罐压等于 $p_s$，则 $\nu_f=p_sV/(R\gamma T_s)$。

若罐壁初温为 $T_{w0}$、热容为 $C_w$，且最终与气体同温，则同一控制体给出
\[
 (\nu_fc_{V,m}+C_w)T_f
 =\nu_fc_{p,m}T_s+C_wT_{w0}.
\]
这一步只改变储能项；气源输入的摩尔焓不变。真实充气过程不可逆，均匀终态可用状态方程描述，充气途中则无须假定罐内始终处于准静态。
\end{solution}


\begin{example}[刚性绝热容器中的余气冷却]\label{ex:cooling}
刚性容器内最初有 $N_0$ 个单原子理想气体粒子，温度 $T_0$。壁绝热，气体经薄壁小孔向真空缓慢泄流，余气始终充分平衡。求余气温度 $T$ 与粒子数 $N$ 的关系。
\Diagram{fig:coolingq}{粒子穿过系统边界并携带能量}{
 \draw[curve](0,0)rectangle(3,2);\fill[white](2.95,.8)rectangle(3.05,1.2);\node at(1.5,1){$N,T,V$};
 \draw[axis](3,1)--(5,1)node[right]{$-\dd N>0$};\node at(4.1,1.65){每粒子 $2\kb T$};\node at(1.5,-.5){刚性、绝热};
}
\end{example}
\begin{solution}
把图\ref{fig:coolingq}中容器内的余气作为变粒子数系统，取 $\dd N<0$。余气能量变化等于流出能量的负值，故
\[
 \dd U=2\kb T\dd N.
\]
另一方面 $U=3N\kb T/2$，全微分必须同时包含 $N$ 与 $T$ 的改变：
\[
 \frac32\kb(N\dd T+T\dd N)=2\kb T\dd N.
\]
移项得到 $3N\dd T=T\dd N$，分离变量并在初末态之间积分，$3\int_{T_0}^{T}\dd\vartheta/\vartheta=\int_{N_0}^{N}\dd N'/N'$。两边的原函数均为对数，代入上下限得
\[
 3\ln\frac T{T_0}=\ln\frac N{N_0},\qquad
 T=T_0\left(\frac N{N_0}\right)^{1/3}.
\]
这里忽略了二阶小量$\dd N\dd T$。逸出粒子比体内平均粒子多携带$\kb T/2$的能量，故余气降温。粒子离开时的$T$应取瞬时余气温度，不能在整个过程中用初温$T_0$计算流出能量。

\Diagram{fig:coolinga}{粒子数减少时温度和压强的变化}{
 \draw[axis](0,0)--(4.8,0)node[right]{$N/N_0$};\draw[axis](0,0)--(0,2.8)node[above]{归一化量};
 \draw[curve]plot[domain=.01:1,samples=60]({4*\x},{2.3*\x^(1/3)});
 \draw[curve,dashed]plot[domain=0:1,samples=60]({4*\x},{2.3*\x^(4/3)});
 \draw[guide](4,0)--(4,2.3)--(0,2.3);\node[below]at(4,0){$1$};\node[left]at(0,2.3){$1$};
 \node at(1.5,2.5){$T/T_0$};\node at(2.8,.7){$p/p_0$};
}
容器体积固定，$p/p_0=NT/(N_0T_0)=(N/N_0)^{4/3}$，如图\ref{fig:coolinga}。这里没有活塞体积功，却有随粒子跨边界的能量流；把“绝热”理解成余气内能不变，会同时漏掉开放系统的能量交换。
\end{solution}

\begin{example}[单原子、双原子和非线性多原子的比较]\label{ex:degrees}
在上一例条件下，分别讨论 $f=3,5,6$，求 $N(T)$、$p(T)$ 和减少一半粒子数后的温度。内部自由度在全过程保持同一激发状态。
\Diagram{fig:degreesq}{三类气体的平动和转动自由度}{
 \foreach \x/\s in{0/单原子,2.6/双原子,5.2/非线性分子}{\node[draw,minimum height=.7cm]at(\x,0){\s};}
 \node at(0,-.8){$f=3$};\node at(2.6,-.8){$f=5$};\node at(5.2,-.8){$f=6$};
}
\end{example}
\begin{solution}
图\ref{fig:degreesq}三类气体分别满足 $U=fN\kb T/2$。流出能量按式\eqref{eq:escapeenergy}取 $(f+1)\kb T/2$，于是
\[
 f(N\dd T+T\dd N)=(f+1)T\dd N
 \quad\Rightarrow\quad \frac{\dd N}{N}=f\frac{\dd T}{T}.
\]
沿用上一例的对数积分，只将系数$3$换为$f$，有$\ln(N/N_0)=f\ln(T/T_0)$，即$N/N_0=(T/T_0)^f$。体积固定，压强还多含一个温度因子，故$p/p_0=(T/T_0)^{f+1}$。
三种情形的内部能量和逸出能量可以按每粒子逐项列出：平动在体内占 $3\kb T/2$，在逸出样本占 $2\kb T$；双原子的两个转动二次项另加 $\kb T$，非线性分子的三个转动项另加 $3\kb T/2$。因此逸出能量依次为 $2\kb T$、$3\kb T$、$7\kb T/2$。穿孔不改变内部状态的平均分布，所以这些转动项在能量守恒中部分抵消，留下的是不同的储热能力与同一个额外平动能。

\Diagram{fig:degreesa}{相同泄流比例下不同自由度的温度下降}{
 \draw[axis](0,0)--(4.7,0)node[right]{$N/N_0$};\draw[axis](0,0)--(0,2.7)node[above]{$T/T_0$};
 \foreach \f in{3,5,6}{\draw[curve]plot[domain=.02:1,samples=50]({4*\x},{2.3*\x^(1/\f)});}
 \draw[guide](4,0)--(4,2.3)--(0,2.3);\node[below]at(4,0){$1$};\node[left]at(0,2.3){$1$};
 \node at(2,2.8){上至下：$f=6,5,3$};
}
当 $N=N_0/2$，三种温度依次为 $T_02^{-1/3}$、$T_02^{-1/5}$、$T_02^{-1/6}$。图\ref{fig:degreesa}中自由度越多，冷却越弱，因为逸出粒子相对于体内平均粒子多带走的能量始终只有 $\kb T/2$，而每个粒子的热容随 $f$ 增大。
\end{solution}

\begin{example}[温度和粒子数的时间规律]\label{ex:coolingtime}
容器体积 $V$、孔面积 $A_h$，其余条件同前。气体自由度 $f$ 恒定。求 $T(t)$、$N(t)$、$p(t)$，以及温度降至 $T_0/2$ 所需时间。
\Diagram{fig:cooltimeq}{泄流速率由孔面积和余气状态决定}{
 \draw(0,0)rectangle(3,1.8);\fill[white](2.95,.7)rectangle(3.05,1.1);\draw[axis](3,.9)--(5,.9);
 \node at(1.5,.9){$V,N(t),T(t)$};\node at(4,1.6){$A_h$};\node at(4,-.1){$\dot N<0$};
}
\end{example}
\begin{solution}
图\ref{fig:cooltimeq}中穿孔粒子数率是碰壁数率乘孔面积：
\[
 \dot N=-\frac{A_hN}{V}\sqrt{\frac{\kb T}{2\pi m}}.
\]
由 $N\propto T^f$，对时间求导 $\dot N/N=f\dot T/T$，故
\[
 \dot T=-\frac{A_h}{fV}\sqrt{\frac{\kb}{2\pi m}}T^{3/2}.
\]
分离变量并取初末态为积分上下限，
\[
 \left[-2\vartheta^{-1/2}\right]_{T_0}^{T}
 =-\frac{A_h}{fV}\sqrt{\frac{\kb}{2\pi m}}\int_0^t\dd s.
\]
代入端点并约去两边负号，得到
\[
 \frac1{\sqrt T}-\frac1{\sqrt{T_0}}
 =\frac{A_h}{2fV}\sqrt{\frac{\kb}{2\pi m}}t.
\]
定义初始数目衰减率 $\kappa_0=(A_h/V)\sqrt{\kb T_0/(2\pi m)}$ 以及 $t_c=2f/\kappa_0$，则
\[
 \frac{T}{T_0}=(1+t/t_c)^{-2},\quad
 \frac{N}{N_0}=(1+t/t_c)^{-2f},\quad
 \frac{p}{p_0}=(1+t/t_c)^{-2(f+1)}.
\]
定义的$\kappa_0$量纲为时间的倒数，是初始时刻的$-\dot N/N$；后续瞬时衰减率为$\kappa_0\sqrt{T/T_0}$，随降温而减小，所以粒子数遵循幂律而非恒温指数律。

\Diagram{fig:cooltimea}{变换后温度的线性时间关系}{
 \draw[axis](0,0)--(5,0)node[right]{$t/t_c$};\draw[axis](0,0)--(0,3.4)node[above]{$\sqrt{T_0/T}$};
 \draw[curve](0,.7)--(4,2.7);\node[left]at(0,.7){$1$};\node at(3,.8){$1+t/t_c$};
}
图\ref{fig:cooltimea}把$\sqrt{T_0/T}=1+t/t_c$画成直线，斜率为$1/t_c$。取$T/T_0=1/2$，有$1+t/t_c=\sqrt2$，故降至一半初温需$t=(\sqrt2-1)t_c$。
\end{solution}

\begin{example}[温变热容下的泄流]
气体单粒子的平衡内能为 $u(T)$，热容 $c_v(T)=\dd u/\dd T$。内部状态与平动独立，穿孔不选择内部状态。求一般的 $N(T)$ 关系，并取 $c_v(T)=c_0+bT$ 求显式结果。
\Diagram{fig:variablecvq}{逸出只对平动速率加权}{
 \node[draw,align=center]at(0,0){体内\\$u(T)$};\draw[axis](1,0)--(3,0);\node[draw,align=center]at(4,0){逸出\\$u(T)+\kb T/2$};
}
\end{example}
\begin{solution}
图\ref{fig:variablecvq}表示内部平均能量不变，只有平动抽样多出 $\kb T/2$。余气总能量为 $Nu(T)$，所以
\[
 Nc_v(T)\dd T+u(T)\dd N=[u(T)+\kb T/2]\dd N.
\]
消去两侧的 $u\dd N$，有 $\dd N/N=2c_v(T)\dd T/(\kb T)$。因此
\[
 \ln\frac N{N_0}=\int_{T_0}^T\frac{2c_v(\vartheta)}{\kb\vartheta}\dd\vartheta.
\]
若给每个分子加上相同的常数参考能，余气和流出粒子的能量同时改变，温度方程不变。温变热容下，单粒子内能差要积$c_v\dd T$，粒子数关系则要积$c_v\dd T/T$；两种权重须分别保留。

\Diagram{fig:variablecva}{温变热容分别贡献对数项和线性项}{
 \node[draw]at(0,0){$c_0/T$};\draw[->](0,-.35)--(0,-1);\node[below]at(0,-1){$c_0\ln(T/T_0)$};
 \node[draw]at(4,0){$b$};\draw[->](4,-.35)--(4,-1);\node[below]at(4,-1){$b(T-T_0)$};
}
按图\ref{fig:variablecva}将积分分成两项，
\[
 \ln\frac{N}{N_0}
 =\frac{2}{\kb}\left[c_0\ln\vartheta+b\vartheta\right]_{T_0}^{T}
 =\frac{2c_0}{\kb}\ln\frac{T}{T_0}+\frac{2b}{\kb}(T-T_0).
\]
两边取指数，得到
\[
 \frac N{N_0}=\left(\frac T{T_0}\right)^{2c_0/\kb}
 \exp\!\left[\frac{2b}{\kb}(T-T_0)\right].
\]
该式只在给定热容模型的适用温区有效；若内部弛豫比泄流还慢，则内部状态不能由同一瞬时温度描述，需要分别写能量交换方程。
\end{solution}

\section{光子气体的压强与辐射}
\label{th:radiation}
热平衡辐射的传播方向各向同性，光子速率为 $c$。若频率为 $\omega/(2\pi)$ 的单光子能量为 $\varepsilon=\hbar\omega$，其动量大小是 $\varepsilon/c$。为避免把粒子数与能量分布混在一起，直接用频率积分后的能量密度 $u$。每个方向的能量密度为 $u\dd\Omega/(4\pi)$。

单位面积、单位时间穿出小孔的能量为
\begin{equation}\label{eq:radiationflux}
 J_E=uc\int_0^{\pi/2}\frac12\sin\theta\cos\theta\dd\theta=\frac14cu.
\end{equation}
若黑体总辐射通量为 $\sigma_{\rm SB}T^4$，便有 $u=aT^4$，$a=4\sigma_{\rm SB}/c$。这里 $u$ 的单位为能量每体积，$J_E$ 为功率每面积；相差的因子 $c/4$ 有明确的速度和几何意义。只靠经典 Maxwell--Boltzmann 分布不能确定辐射谱和常数 $a$，本讲不把光子当作守恒的经典分子。

\begin{example}[腔内能量与小孔辐射]\label{ex:cavity}
温度 $T$ 的黑体腔体积为 $V$，小孔面积为 $A_h$。已知 $\sigma_{\rm SB}$，求腔内辐射能量与经小孔向冷真空辐射的功率。若外界也是温度 $T_e$ 的黑体辐射场，净功率如何改变？
\Diagram{fig:cavityq}{小孔连接腔内与外界辐射场}{
 \draw(0,0)rectangle(3,1.8);\fill[white](2.95,.7)rectangle(3.05,1.1);\node at(1.5,.9){$T,V$};
 \draw[axis](3,.9)--(5,.9)node[right]{$T_e$};\node at(4,1.5){$A_h$};
}
\end{example}
\begin{solution}
图\ref{fig:cavityq}中腔内能量密度为 $u=4\sigma_{\rm SB}T^4/c$，故 $U=uV=4\sigma_{\rm SB}VT^4/c$。由式\eqref{eq:radiationflux}，出射功率
\[
 P_{\rm out}=A_h\frac{cu}{4}=A_h\sigma_{\rm SB}T^4.
\]
小孔的出射通量由孔口处的各向同性辐射场决定，乘的是孔面积 $A_h$。腔体其余内壁虽然也吸收、发射光，但它们之间的交换留在腔内，不应把整个腔壁面积再乘进外流功率。外界有平衡辐射场时，两个单向功率分别存在，温度相同只意味着净值为零。$U$ 的量纲是焦耳，$P$ 的量纲是瓦特，连接两者的面积、体积和速度因子可由 $P/U=A_hc/(4V)$ 直接辨认。

\Diagram{fig:cavitya}{两侧能流分别计算后相减}{
 \draw[axis](0,.5)--(4,.5)node[midway,above]{$A_h\sigma_{\rm SB}T^4$};
 \draw[axis](4,-.5)--(0,-.5)node[midway,below]{$A_h\sigma_{\rm SB}T_e^4$};
}
图\ref{fig:cavitya}中，外界各向同性辐射的入射功率为 $A_h\sigma_{\rm SB}T_e^4$，相减得净外流功率 $P=A_h\sigma_{\rm SB}(T^4-T_e^4)$。该式要求腔内辐射始终接近平衡，且孔口接收完整的入射半球。孔过大导致明显不均匀，或孔前存在方向筛选时，须根据实际辐射分布重新作角积分。
\end{solution}

\begin{example}[辐射压强]\label{ex:photonpressure}
能量密度为 $u$ 的各向同性辐射照射理想反射平面。由光子动量变化求壁面压强。
\Diagram{fig:photonpressureq}{反射光子的法向动量变化}{
 \draw[line width=1.3pt](3,-1.2)--(3,1.5);\draw[axis](.3,1.2)--(3,0);\draw[axis](3,0)--(.3,-1.2);\draw[guide](1,0)--(4,0);
 \node at(1.4,.9){$\varepsilon/c$};\node at(2,.3){$\theta$};
}
\end{example}
\begin{solution}
图\ref{fig:photonpressureq}的单光子法向冲量为 $2\varepsilon\cos\theta/c$。把方向能量密度 $u\dd\Omega/(4\pi)$ 除以单光子能量，再乘扫掠速度 $c\cos\theta$ 与冲量，能量和速度中的因子分别约去：
\[
 p=\int_{\rm hem}\frac u\varepsilon\frac{\dd\Omega}{4\pi}
 c\cos\theta\frac{2\varepsilon}c\cos\theta
 =u\int_0^{\pi/2}\sin\theta\cos^2\theta\dd\theta
 =u\left[-\frac{\cos^3\theta}{3}\right]_0^{\pi/2}=\frac u3.
\]
频率分布不必预先给定。若用单位频率能量密度 $u_\omega$ 计算，每个频段的动量流同样给出 $p_\omega=u_\omega/3$，再对频率积分即可。因此 $p=u/3$ 依赖的是光子的能量动量关系和角度各向同性，不依赖这里是否知道具体的黑体谱。法向入射的单向光束则没有半球角度平均；对同样的入射能流，吸收与镜面反射的压强还相差两倍。

\Diagram{fig:photonpressurea}{角度平均将动量流分成三个等价方向}{
 \draw[axis](0,0)--(2,0)node[right]{$x$};\draw[axis](0,0)--(0,1.8)node[above]{$z$};\draw[axis](0,0)--(-1.4,-.9)node[left]{$y$};\node at(3,1){$\langle\cos^2\theta\rangle_{4\pi}=1/3$};
}
图\ref{fig:photonpressurea}给出同一结果的对称性解释。对黑色吸收壁，入射光的压强只有上述反射结果的一半；处于辐射平衡的壁还会向腔内发射光，发射反冲贡献另一半，总压强仍为 $u/3$。单向光束则不满足各向同性假设。
\end{solution}

\section{热力学势与偏导关系}
\label{th:potentials}
对成分、物质的量固定的简单可压缩系统，平衡态基本关系为
\begin{equation}\label{eq:fundamental}
 \dd U=T\dd S-p\dd V.
\end{equation}
这是状态空间中的微分恒等式，不要求实际过程可逆；但只有可逆传热才可把 $T\dd S$ 认作过程热量 $\delta Q_{\rm rev}$。不可逆过程中一般 $\delta Q\ne T\dd S$，边界做功也未必是内部平衡压强乘 $\dd V$。

以 $U(S,V)$ 为势，$T=(\partial U/\partial S)_V$、$p=-(\partial U/\partial V)_S$。等温条件下希望把 $T$ 作为自变量，便减去 $TS$；等压条件下希望把 $p$ 作为自变量，便加上 $pV$。于是四种势为
\[
 \begin{array}{c|c|c}
 \text{势}&\text{自然变量}&\text{全微分}\\\hline
 U&S,V&T\dd S-p\dd V\\
 F=U-TS&T,V&-S\dd T-p\dd V\\
 H=U+pV&S,p&T\dd S+V\dd p\\
 G=U+pV-TS&T,p&-S\dd T+V\dd p
 \end{array}
\]
对同一势的混合二阶导数交换次序，依次得到
\begin{align}
 -\left(\frac{\partial p}{\partial S}\right)_V
 &=\left(\frac{\partial T}{\partial V}\right)_S,&
 \left(\frac{\partial S}{\partial V}\right)_T
 &=\left(\frac{\partial p}{\partial T}\right)_V,\\
 \left(\frac{\partial T}{\partial p}\right)_S
 &=\left(\frac{\partial V}{\partial S}\right)_p,&
 \left(\frac{\partial S}{\partial p}\right)_T
 &=-\left(\frac{\partial V}{\partial T}\right)_p.
\end{align}
记忆关系式时保留自然变量，比仅记交换字母可靠。状态方程 $p=p(V,T)$ 可以给出压强的偏导，却不能独自确定任意温度下的内能；还需一个热容函数以及参考值，并满足相容条件。

\subsection{由状态方程求内能}
\label{th:energyidentity}
从 $U(V,S)$ 改用 $V,T$ 作坐标，需要沿等温方向改变 $V$ 时同时计入 $S(V,T)$ 的改变：
\begin{align}
 \left(\frac{\partial U}{\partial V}\right)_T
 &=\left(\frac{\partial U}{\partial V}\right)_S
 +\left(\frac{\partial U}{\partial S}\right)_V
 \left(\frac{\partial S}{\partial V}\right)_T\\
 &=-p+T\left(\frac{\partial p}{\partial T}\right)_V.
\end{align}
又定义 $C_V=(\partial U/\partial T)_V$，所以
\begin{equation}\label{eq:duvt}
 \dd U=C_V\dd T+\left[T\left(\frac{\partial p}{\partial T}\right)_V-p\right]\dd V.
\end{equation}
代入式\eqref{eq:fundamental}，得到计算熵的常用式
\begin{equation}\label{eq:dsvt}
 \dd S=\frac{C_V}{T}\dd T+
 \left(\frac{\partial p}{\partial T}\right)_V\dd V.
\end{equation}
由于$U$是状态函数，$\partial_VU_T=\partial_TU_V$。右侧的乘积求导给出$\partial_T(Tp_T-p)=p_T+Tp_{TT}-p_T=Tp_{TT}$，故
\begin{equation}\label{eq:cvcompat}
 \left(\frac{\partial C_V}{\partial V}\right)_T
 =T\left(\frac{\partial^2p}{\partial T^2}\right)_V.
\end{equation}
因此指定状态方程以后，不能任意再指定一个与它冲突的热容。

\subsection{焓与定压热容}
\label{th:enthalpy}
由$\dd H=T\dd S+V\dd p$，在定压可逆过程中$\delta Q=\dd H$，故$C_p=(\partial H/\partial T)_p$。写$\dd S=S_T|_p\dd T+S_p|_T\dd p$，其中$S_T|_p=C_p/T$；由$G(T,p)$的混合偏导有$S_p|_T=-V_T|_p$。代回焓微分，得到
\begin{equation}\label{eq:dhpt}
 \dd H=C_p\dd T+
 \left[V-T\left(\frac{\partial V}{\partial T}\right)_p\right]\dd p,
 \qquad
 \dd S=\frac{C_p}{T}\dd T-
 \left(\frac{\partial V}{\partial T}\right)_p\dd p.
\end{equation}
在单相稳定区，定义热膨胀系数 $\alpha_p=V^{-1}(\partial V/\partial T)_p$、等温压缩率 $\kappa_T=-V^{-1}(\partial V/\partial p)_T$，可将热容差写为
\begin{equation}\label{eq:cpminuscv}
 C_p-C_V=T\left(\frac{\partial p}{\partial T}\right)_V
 \left(\frac{\partial V}{\partial T}\right)_p
 =\frac{TV\alpha_p^2}{\kappa_T}.
\end{equation}
推导时在式\eqref{eq:dsvt}中沿定压方向取 $\dd V=(\partial V/\partial T)_p\dd T$，再与定压熵微分比较即可。稳定性要求 $\kappa_T>0$，从而 $C_p\ge C_V$。

\subsection{自然变量怎样进入偏导数}
\label{th:naturalderivatives}
势函数的自然变量来自其全微分中独立微分的系数。以 $F=U-TS$ 为例，展开乘积，
\[
 \dd F=\dd U-T\dd S-S\dd T
 =-S\dd T-p\dd V.
\]
$T\dd S$ 的抵消使 $F$ 可直接用 $T,V$ 作为自变量。于是
$S=-F_T|_V$、$p=-F_V|_T$；再求混合导数，
\[
 \left(\frac{\partial S}{\partial V}\right)_T
 =-F_{TV}
 =-F_{VT}
 =\left(\frac{\partial p}{\partial T}\right)_V.
\]
等式中两个负号来自同一个势的两个一阶导数。若从 $U$ 或 $H$ 出发，系数中的正负号不同，所得 Maxwell 关系也随之改变。实际使用时，先写势的全微分，再读出一阶导数，可直接确定符号。

“自然变量”不表示势只能在这些变量下表示。$U(S,V)$ 可改写成 $U(T,V)$，只是后者的偏导数须经过链式法则。沿 $V$ 固定的升温方向，
\[
 \left(\frac{\partial U}{\partial T}\right)_V
 =T\left(\frac{\partial S}{\partial T}\right)_V=C_V.
\]
沿 $T$ 固定的膨胀方向，熵一般同时改变，才有式\eqref{eq:duvt}中的体积系数。若把两条不同方向上的偏导数视为同一个“内能对体积的变化率”，便会丢掉 $T S_V|_T$ 这一项。

\Diagram{fig:derivativepaths}{同一点出发的等温方向与等熵方向}{
 \draw[axis](0,0)--(5.5,0)node[right]{$V$};
 \draw[axis](0,0)--(0,3.2)node[above]{$T$};
 \coordinate(A)at(1,2.4);\fill(A)circle(1.4pt)node[above left]{$A$};
 \draw[axis](A)--(4.6,2.4)node[right]{$T$ 固定};
 \draw[curve]plot[domain=1:4.6](\x,{2.4/\x^.45});
 \draw[->](2.1,1.72)--(2.6,1.56);
 \node at(3.8,.75){$S$ 固定};
}
图\ref{fig:derivativepaths}示意理想气体的两条方向。在等熵方向，内能变化可直接用 $-p\dd V$；沿等温方向则必须同时计算熵变。对于理想气体，两项恰好抵消，使等温内能不变；这一抵消是具体状态方程的结果。

\subsection{全微分、倒数关系与循环关系}
设单相区的状态方程为 $p=p(T,V)$，全微分为
$\dd p=p_T|_V\dd T+p_V|_T\dd V$。沿定压方向令 $\dd p=0$，得到
\[
 \left(\frac{\partial V}{\partial T}\right)_p
 =-\frac{(\partial p/\partial T)_V}
 {(\partial p/\partial V)_T}.
\]
这个比值是状态面上定压切线的斜率。保持同一个第三变量不变时，两个互为反函数的局部斜率满足倒数关系；保持条件不同时，则须从全微分重新推导。将下列三个偏导相乘，得到循环关系
\[
 \left(\frac{\partial p}{\partial T}\right)_V
 \left(\frac{\partial T}{\partial V}\right)_p
 \left(\frac{\partial V}{\partial p}\right)_T=-1.
\]
循环关系中的负号来自 $\dd p=0$ 时两项相抵。若某个分母偏导为零或发散，所用坐标可能失去局部可逆性，应回到全微分或改用另一组状态变量，而不能继续机械使用倒数。

\subsection{热容差、压缩率与绝热斜率}
\label{th:compressibilityratio}
记 $A=(\partial p/\partial T)_V$、$B=(\partial p/\partial V)_T$，其中稳定单相区有 $B<0$。沿定压方向，$\dd V/\dd T=-A/B$，代入熵微分，
\[
 C_p=C_V-\frac{TA^2}{B}.
\]
该式与响应系数形式的式\eqref{eq:cpminuscv}相同。沿等熵方向，先由式\eqref{eq:dsvt}求
\[
 \left(\frac{\partial T}{\partial V}\right)_S
 =-\frac{TA}{C_V},
\]
再代入状态方程全微分，得到
\[
 \left(\frac{\partial p}{\partial V}\right)_S
 =B-\frac{TA^2}{C_V}
 =B\,\frac{C_p}{C_V}.
\]
定义 $\kappa_S=-V^{-1}(\partial V/\partial p)_S$，便有
\[
 \frac{\kappa_T}{\kappa_S}=\frac{C_p}{C_V}.
\]
在 $C_V>0,\kappa_T>0$ 的稳定区，绝热压缩比等温压缩更困难，因为压缩引起的升温会进一步增大压强。理想气体给 $\kappa_T=1/p$、$\kappa_S=1/(\gamma p)$；实际气体也满足上述热容比关系，但热容比一般随状态变化。

小振幅声波中的一小团气体在振动周期内来不及与周围交换足够热量时，其压强密度关系近似等熵。由于 $\rho\propto1/V$，
\[
 c_s^2=\left(\frac{\partial p}{\partial\rho}\right)_S
 =-\frac{V}{\rho}\left(\frac{\partial p}{\partial V}\right)_S
 =\frac1{\rho\kappa_S}.
\]
理想气体因而有 $c_s^2=\gamma RT/M$。这里忽略了黏滞、热传导和内部自由度弛豫对声速的频率依赖；采用哪一热容，应与声波时间尺度内能否建立内部平衡相对应。

\subsection{实际气体中两种压缩斜率的比较}
用摩尔体积 $v$ 表示 van der Waals 状态方程
$p=RT/(v-b)-a/v^2$，摩尔定容热容为常数 $c$。在一摩尔的熵微分中，
$s_T|_v=c/T$、$s_v|_T=R/(v-b)$，于是等熵方向的温度变化为
\[
 \left(\frac{\partial T}{\partial v}\right)_s
 =-\frac{RT}{c(v-b)}.
\]
对状态方程作链式求导，可得
\begin{align*}
 \left(\frac{\partial p}{\partial v}\right)_s
 &=\left(\frac{\partial p}{\partial v}\right)_T+
 \left(\frac{\partial p}{\partial T}\right)_v
 \left(\frac{\partial T}{\partial v}\right)_s\\
 &=-\left(1+\frac Rc\right)\frac{RT}{(v-b)^2}
 +\frac{2a}{v^3}.
\end{align*}
令 $\gamma_0=1+R/c$，便有
\[
 c_s^2=\frac1M\left[
 \frac{\gamma_0RTv^2}{(v-b)^2}-\frac{2a}{v}\right].
\]
这里先完成等熵导数，再乘 $\dd v/\dd\rho=-v^2/M$。$\gamma_0$ 是这个常定容热容模型中出现的组合，实际的 $C_p/C_V$ 应由前文热容差求出，通常与 $\gamma_0$ 不同。

第一项给出热运动及排斥体积对恢复力的贡献，吸引项则削弱这一响应。令 $a,b\to0$，结果恢复为理想气体声速。实际气体的声速由状态方程的等熵斜率确定，既取决于压强，也取决于它随密度和温度变化的方式。

还可在等温失稳边界附近直接比较斜率。若 $p_v|_T\to0^-$，由前面的通式有
$p_v|_s\to-TA^2/C_V$，在 $A$ 与 $C_V$ 有限时仍可为有限负值。等温压缩率发散与等熵响应发散因而不是同一个条件。这一比较限于给定状态方程及常热容近似；接近真实临界点时，热容和弛豫时间本身也可能显著变化。


\begin{example}[理想气体的内能与熵]\label{ex:idealstate}
$\nu$ 摩尔理想气体满足 $pV=\nu RT$，已知 $C_V$ 为常量。相对于参考态 $(V_0,T_0)$，求 $U$、$H$、$S$ 的变化，说明为什么内能与体积无关。
\Diagram{fig:idealstateq}{参考态到任意态的两段积分路径}{
 \draw[axis](0,0)--(4.8,0)node[right]{$V$};\draw[axis](0,0)--(0,3)node[above]{$T$};
 \coordinate(a)at(1,.8);\coordinate(b)at(3.5,.8);\coordinate(c)at(3.5,2.3);
 \draw[axis](a)--(b);\draw[axis](b)--(c);\node[below]at(a){$A(V_0,T_0)$};\node[right]at(c){$C(V,T)$};\node[below]at(b){$B$};
}
\end{example}
\begin{solution}
$p_T|_V=\nu R/V$，所以$U_V|_T=T\nu R/V-p=0$。图\ref{fig:idealstateq}的等温段$A\to B$不改变内能；等容段$B\to C$给出$U-U_0=\int_{T_0}^TC_V\dd\vartheta=[C_V\vartheta]_{T_0}^T=C_V(T-T_0)$。焓增量为
\[
 H-H_0=U-U_0+pV-p_0V_0=(C_V+\nu R)(T-T_0),
\]
故 $C_p=C_V+\nu R$。熵微分为 $\dd S=C_V\dd T/T+\nu R\dd V/V$。
图中第一段固定 $T_0$、第二段固定最终体积 $V$，是为积分选取的可逆参考路径，并不限定系统实际如何由初态到末态。两段内能变化相加时，等温段为零，等容段才贡献热容积分。焓中还含 $pV=\nu RT$，所以定压热容多出 $\nu R$；这里 $C_V,C_p$ 均为整个 $\nu$ 摩尔系统的热容。把它们换成摩尔量时，内能与熵的温度项还要同时乘回物质的量。

\Diagram{fig:idealstatea}{两段路径分别给出体积项与温度项}{
 \node[draw,align=center]at(0,0){$A\to B$ 等温\\$\nu R\ln(V/V_0)$};
 \node[draw,align=center]at(4.4,0){$B\to C$ 等容\\$C_V\ln(T/T_0)$};\draw[->](1.5,-.8)--(3,-.8);
}
按图\ref{fig:idealstatea}，两段熵变分别为$\nu R\int_{V_0}^V\dd w/w=\nu R[\ln w]_{V_0}^V$和$C_V\int_{T_0}^T\dd\vartheta/\vartheta=C_V[\ln\vartheta]_{T_0}^T$，相加得
\[
 S-S_0=C_V\ln\frac T{T_0}+\nu R\ln\frac V{V_0}.
\]
内能仅依赖温度并非“所有气体”的性质，而是理想气体状态方程与热力学相容性的结果。若 $C_V=C_V(T)$，内能温度项改为 $\int C_V\dd T$，熵温度项改为 $\int C_V\dd T/T$。
\end{solution}

\begin{example}[van der Waals 气体的内能与熵]\label{ex:vdwstate}
固定物质的量 $\nu$，气体满足
\[
 p=\frac{\nu RT}{V-B}-\frac A{V^2},\quad A=\nu^2a,\ B=\nu b,
\]
其中 $a,b$ 为摩尔形式的常数。给定常量 $C_V$，求相对于参考态 $(V_0,T_0)$ 的内能和熵变化。只讨论稳定单相区，$V>B$。
\Diagram{fig:vdwstateq}{排除体积改变可用的构型空间}{
 \draw[axis](0,0)--(5,0)node[right]{$V$};\draw[axis](0,0)--(0,2.3)node[above]{$p$};
 \draw[guide](1,0)--(1,2.1)node[above]{$B$};\draw[curve]plot[domain=1.6:4.6,samples=50](\x,{1/(\x-1)-.5/\x^2});\node at(3.6,1.8){稳定区示意};
}
\end{example}
\begin{solution}
图\ref{fig:vdwstateq}的 $B$ 是固定量，求偏导时不随 $T$ 改变。有 $p_T|_V=\nu R/(V-B)$，代入式\eqref{eq:duvt}：
\[
 T p_T-p=\frac A{V^2},\qquad
 \dd U=C_V\dd T+\frac A{V^2}\dd V.
\]
分别积分温度项和体积项。温度积分同上一例；体积积分为$\int_{V_0}^{V}A\dd w/w^2=[-A/w]_{V_0}^{V}=A(1/V_0-1/V)$，所以
\[
 U-U_0=C_V(T-T_0)+A\left(\frac1{V_0}-\frac1V\right).
\]
这里$A,B$均按固定物质的量看作常数。原函数$-A/V$的导数为$+A/V^2$，因此膨胀使负的吸引势能升高。参考常数只改变内能与熵的零点，不改变偏导、热容或两态之间的差。

\Diagram{fig:vdwstatea}{熵微分中的吸引项相消}{
 \node[draw,align=center]at(0,0){$\dd U$ 中\\$+A\dd V/V^2$};\node at(2.7,0){$+$};
 \node[draw,align=center]at(5.1,0){$p\dd V$ 中\\$-A\dd V/V^2$};\node at(2.5,-1){$T\dd S=C_V\dd T+\nu RT\dd V/(V-B)$};
}
按图\ref{fig:vdwstatea}在$\dd U+p\dd V$中消去吸引项，再除以$T$。温度项仍给$C_V\ln(T/T_0)$；体积项令$z=w-B$、$\dd z=\dd w$，上下限变为$V_0-B,V-B$，于是$\nu R\int_{V_0-B}^{V-B}\dd z/z=\nu R[\ln z]_{V_0-B}^{V-B}$。两项相加，
\[
 S-S_0=C_V\ln\frac T{T_0}+\nu R\ln\frac{V-B}{V_0-B}.
\]
吸引参数 $A$ 出现在内能中，但在以 $V,T$ 为自变量的熵差中抵消；排除体积 $B$ 则保留。若把 $A,B$ 当成摩尔常数却仍用总容积 $V$，会漏掉物质的量因子。
\end{solution}

\begin{example}[实际气体的可逆绝热膨胀]\label{ex:vdwadiabat}
上述 van der Waals 气体由 $(V_1,T_1)$ 可逆绝热膨胀至体积 $V_2$，热容为常量。求 $T_2$、$p_2$ 和气体对外做功，并比较同初温、同体积比的理想气体。
\Diagram{fig:vdwadiabatq}{绝热线上的初末两态}{
 \draw[axis](0,0)--(5,0)node[right]{$V$};\draw[axis](0,0)--(0,3)node[above]{$p$};
 \draw[curve]plot[domain=1:4.3](\x,{1.8/(\x-.3)^1.3-.2/\x^2});\fill(1,{1.8/.7^1.3-.2})circle(1.2pt)node[above right]{$1$};\fill(4,{1.8/3.7^1.3-.2/16})circle(1.2pt)node[above right]{$2$};\draw[->]plot[domain=2:2.5](\x,{1.8/(\x-.3)^1.3-.2/\x^2});
}
\end{example}
\begin{solution}
图\ref{fig:vdwadiabatq}的路径可逆绝热，故 $\dd S=0$。由上一例的熵微分
\[
 C_V\frac{\dd T}T+\nu R\frac{\dd V}{V-B}=0.
\]
取初态$1$为积分下限，利用上一例的两个对数积分，得到$C_V\ln(T/T_1)+\nu R\ln[(V-B)/(V_1-B)]=0$，即$T(V-B)^{\nu R/C_V}=T_1(V_1-B)^{\nu R/C_V}$。
所以 $T_2=T_1[(V_1-B)/(V_2-B)]^{\nu R/C_V}$，压强再由状态方程求得
\[
 p_2=\frac{\nu RT_2}{V_2-B}-\frac A{V_2^2}.
\]
令 $s=\nu R/C_V$，将绝热温度关系代入热压强项，可写成 $T(V)=K(V-B)^{-s}$。其功积分中的热运动部分为
\begin{align*}
 \nu RK\int_{V_1}^{V_2}(V-B)^{-s-1}\dd V
 &=\left[-\frac{\nu RK}{s}(V-B)^{-s}\right]_{V_1}^{V_2}\\
 &=\frac{\nu RK}{s}\bigl[(V_1-B)^{-s}-(V_2-B)^{-s}\bigr].
\end{align*}
因$\nu R/s=C_V$，热运动部分为$C_V(T_1-T_2)$。吸引项单独积分为$\int_{V_1}^{V_2}(-A/V^2)\dd V=[A/V]_{V_1}^{V_2}=-A(1/V_1-1/V_2)$，代表克服分子间吸引所需的能量。两项之和即气体对外做的功。

\Diagram{fig:vdwadiabata}{绝热功等于内能减少}{
 \node[draw]at(0,0){热运动能减少};\node[draw]at(4,0){吸引势能增加};\draw[->](0,-.4)--(2,-1);\draw[->](4,-.4)--(2,-1);\node[below]at(2,-1){$W=C_V(T_1-T_2)-A(1/V_1-1/V_2)$};
}
第一定律给 $W=U_1-U_2$，即图\ref{fig:vdwadiabata}的表达式。吸引势能在膨胀时上升，因而减少可供做功的热运动能。对于相同 $V_2/V_1>1$，$B>0$ 使 $(V_1-B)/(V_2-B)<V_1/V_2$，温度比理想气体降得更多；但功还包含吸引能，不能只比较温降。等价的绝热线为 $(p+A/V^2)(V-B)^{1+\nu R/C_V}=\text{常量}$。
\end{solution}

\begin{example}[从辐射压强求能量与熵]\label{ex:photonstate}
平衡光子气体满足 $p=b_rT^4$，$b_r$ 是常量。利用热力学关系求 $U(V,T)$、$S(V,T)$。说明不守恒的光子数为何不妨碍使用 $\dd U=T\dd S-p\dd V$。
\Diagram{fig:photonstateq}{可改变光子数的平衡辐射腔}{
 \draw(0,0)rectangle(3.5,1.8);\draw[decorate,decoration={snake,amplitude=1.3pt,segment length=7pt},->](.4,.5)--(3,.5);
 \draw[decorate,decoration={snake,amplitude=1.3pt,segment length=7pt},->](3,1.3)--(.4,1.3);\node[right]at(3.7,.9){$p=b_rT^4$};
}
\end{example}
\begin{solution}
一般开放系统的基本关系中还含 $\mu\dd N$。图\ref{fig:photonstateq}的热平衡光子可以被壁吸收或发射，光子化学势为零，因此该项消失。用式\eqref{eq:duvt}，
\[
 \left(\frac{\partial U}{\partial V}\right)_T=-b_rT^4+4b_rT^4=3b_rT^4.
\]
先积分得 $U=3b_rVT^4+\phi(T)$。均匀系统在固定温度下能量随体积广延，$U(\lambda V,T)=\lambda U(V,T)$，故不能保留体积无关的 $\phi(T)$。于是 $U=3b_rVT^4$。
广延性也可通过复制腔体理解：同温、同体积的两份辐射合在一起，忽略界面效应，能量应加倍。求出下文的熵后，代入$G=U+pV-TS$得到$G=0$，与平衡光子的零化学势一致。固定粒子数的分子气体仍须保留其粒子数约束。

\Diagram{fig:photonstatea}{在状态平面中积分熵微分}{
 \draw[axis](0,0)--(4.5,0)node[right]{$V$};\draw[axis](0,0)--(0,2.8)node[above]{$T$};
 \draw[axis](1,.7)--(3.2,.7);\draw[axis](3.2,.7)--(3.2,2.1);\node[below]at(1,.7){$(V_0,T_0)$};\node[right]at(3.2,2.1){$(V,T)$};
}
对$U=3b_rVT^4$作乘积微分，$\dd U=3b_rT^4\dd V+12b_rVT^3\dd T$。加$p\dd V$并除以$T$，得$\dd S=4b_rT^3\dd V+12b_rVT^2\dd T$。按图\ref{fig:photonstatea}先沿$T_0$等温段积分，再沿$V$等容段积分，
\begin{align*}
 \Delta S_{\rm iso}&=\int_{V_0}^{V}4b_rT_0^3\dd w=4b_rT_0^3(V-V_0),\\
 \Delta S_{\rm vol}&=\int_{T_0}^{T}12b_rV\vartheta^2\dd\vartheta
 =[4b_rV\vartheta^3]_{T_0}^{T}=4b_rV(T^3-T_0^3).
\end{align*}
相加并消去$VT_0^3$项，得到
\[
 S-S_0=4b_r(VT^3-V_0T_0^3).
\]
取无辐射的零温极限为熵零点，则 $S=4b_rVT^3$。与 $u=aT^4$ 对比有 $a=3b_r$，从而 $S=4aVT^3/3$。
\end{solution}

\begin{example}[状态方程与热容的相容性]
某系统状态方程为 $p=\nu RT/V+qT^2/V^2$，其中 $q>0$。有人同时给定 $C_V=C_0$，这是否相容？若规定在参考体积 $V_*$ 处 $C_V=C_*(T)$，求相容的 $C_V(V,T)$ 与内能。
\Diagram{fig:compatq}{在不同体积处比较定容升温}{
 \draw[axis](0,0)--(4.5,0)node[right]{$V$};\draw[axis](0,0)--(0,2.7)node[above]{$T$};
 \draw[axis](1,.5)--(1,2.1);\draw[axis](3,.5)--(3,2.1);\node[below]at(1,0){$V_*$};\node[below]at(3,0){$V$};
}
\end{example}
\begin{solution}
$p_{TT}|_V=2q/V^2$，故相容式\eqref{eq:cvcompat}要求 $(C_V)_V=2qT/V^2$，不能为零。图\ref{fig:compatq}中两条定容路径的热容存在体积差。积分得
\[
 C_V(V,T)=C_*(T)+2qT\left(\frac1{V_*}-\frac1V\right).
\]
这一结果来自相容条件的体积积分：
\[
 C_V(V,T)-C_V(V_*,T)
 =\int_{V_*}^{V}\frac{2qT}{w^2}\dd w
 =\left[-\frac{2qT}{w}\right]_{V_*}^{V}
 =2qT\left(\frac1{V_*}-\frac1V\right).
\]
再由$U_V=Tp_T-p=qT^2/V^2$积分，得$U=\Phi(T)-qT^2/V$。其温度导数为$\Phi'(T)-2qT/V$；在$V=V_*$匹配热容，确定$\Phi'=C_*(T)+2qT/V_*$。参考热容只规定一条等容线，其余体积处的热容由状态相容性确定。

\Diagram{fig:compata}{同一状态函数的两条积分路径必须闭合}{
 \draw[axis](0,0)--(3,0);\draw[axis](3,0)--(3,1.7);\draw[axis](3,1.7)--(0,1.7);\draw[axis](0,1.7)--(0,0);\node at(1.5,.85){$\oint\dd U=0$};
}
相对于$(V_*,T_*)$，先在$V_*$上从$T_*$升温至$T$，再在终温下从$V_*$改变体积到$V$。第二段的内能增量为$\int_{V_*}^VqT^2w^{-2}\dd w=[-qT^2/w]_{V_*}^V$，所以
\[
 U-U_*=\int_{T_*}^{T}C_*(\vartheta)\dd\vartheta
 +qT^2\left(\frac1{V_*}-\frac1V\right).
\]
图\ref{fig:compata}中，内能沿闭合矩形的增量应为零。若另行指定与状态方程不相容的常热容，两条路径将给出不同内能，模型便失去唯一的状态函数。
\end{solution}

\subsection{由状态关系到过程方程}
\label{th:processclosure}
状态方程规定允许的平衡态，过程约束规定系统怎样在这些状态之间变化。对固定物质的量，式\eqref{eq:duvt}给出内能微分，再用第一定律才能确定热量和功。若过程准静态且只有体积功，$\delta W=p\dd V$，于是
\[
 \delta Q=C_V\dd T+
 T\left(\frac{\partial p}{\partial T}\right)_V\dd V.
\]
这一式与气体的$T\dd S$相同；若外界热源与气体存在有限温差，气体和热源组成的整体仍会产生熵，不能仅凭此式断定整体过程可逆。若外压与内部压强不同，实际功应按边界的外压计算，差值须在不可逆能量转化和熵产生中处理。

给定过程 $V=V(T)$ 后，沿这一路径的热量与温变之比为
\[
 C_{\cal P}=\frac{\delta Q}{\dd T}
 =C_V+T\left(\frac{\partial p}{\partial T}\right)_V
 \frac{\dd V}{\dd T}.
\]
它取决于路径，通常既不是$C_V$也不是$C_p$。路径热容可以为负，即吸热时温度下降。例如对内能只随温度变化且$C_V>0$的理想气体，若吸收热量少于同时对外做的功，差额由内能支付，温度便下降。实际气体的内能还可能含体积项，须联立完整的$\dd U$判断温变。平衡稳定性中的热容条件针对指定约束下的状态响应，应明确保持哪些变量不变。

\begin{example}[多方过程的热量、功与路径热容]
\label{ex:polytropicheat}
$\nu$ 摩尔常热容理想气体经历可逆过程 $pV^a=\text{常量}$，初末体积为 $V_1,V_2$。已知 $C_V=\nu c_{V,m}$、$\gamma=1+\nu R/C_V$。求 $a\ne1$ 时的功、热量和路径热容，并说明 $1<a<\gamma$ 的膨胀过程为什么吸热却降温。
\Diagram{fig:polytropicheatq}{从同一初态出发的等温、多方与绝热线}{
 \draw[axis](0,0)--(5,0)node[right]{$V/V_1$};
 \draw[axis](0,0)--(0,3.2)node[above]{$p/p_1$};
 \foreach \a/\s in{1/solid,1.2/dashed,1.4/densely dotted}
  \draw[curve,\s]plot[domain=1:4.5](\x,{2.5/\x^\a});
 \fill(1,2.5)circle(1.2pt)node[above]{$A$};
 \node at(3.7,2.4){实线 $a=1$};
 \node at(3.7,1.9){虚线 $a=1.2$};
 \node at(3.7,1.4){点线 $\gamma=1.4$};
}
\end{example}
\begin{solution}
图\ref{fig:polytropicheatq}的三条曲线经过同一初态。由理想气体状态方程，
$TV^{a-1}=T_1V_1^{a-1}$，故
$T_2/T_1=(V_1/V_2)^{a-1}$。令 $K=p_1V_1^a$，逐项积分得
\[
 W=\int_{V_1}^{V_2}KV^{-a}\dd V
 =\left[\frac{KV^{1-a}}{1-a}\right]_{V_1}^{V_2}
 =\frac{p_2V_2-p_1V_1}{1-a}
 =\frac{\nu R(T_2-T_1)}{1-a}.
\]
再由 $\Delta U=C_V(T_2-T_1)$ 和 $Q=\Delta U+W$，
\[
 Q=C_V\frac{\gamma-a}{1-a}(T_2-T_1),
 \qquad C_{\cal P}=C_V\frac{\gamma-a}{1-a}.
\]
\Diagram{fig:polytropicheata}{路径热容的符号由两个特征指数分隔}{
 \draw[axis](0,0)--(6.7,0)node[right]{$a$};
 \draw(2.2,-.15)--(2.2,.15)node[above]{$1$};
 \draw(4.4,-.15)--(4.4,.15)node[above]{$\gamma$};
 \node at(1,-.6){$C_{\cal P}>0$};
 \node at(3.3,-.6){$C_{\cal P}<0$};
 \node at(5.5,-.6){$C_{\cal P}>0$};
 \node at(3.3,1){$a=1$ 等温；$a=\gamma$ 可逆绝热};
}
在图\ref{fig:polytropicheata}的中间区间，膨胀给$T_2<T_1$，而$C_{\cal P}<0$，故$Q>0$。功仍为正，且$W>Q$，多出的功来自内能减少。定压过程$a=0$恢复$C_{\cal P}=C_p$；$a=\gamma$给$Q=0$。$a=1$时温变为零，应单独积分：$W=\nu RT\int_{V_1}^{V_2}\dd V/V=\nu RT[\ln V]_{V_1}^{V_2}=\nu RT\ln(V_2/V_1)$，且$\Delta U=0$给$Q=W$，此时热量除以温变不定义有限的路径热容。
\end{solution}

\subsection{内能与熵积分中的参考函数}
由一个偏导积分得到的“常数”可能仍依赖其他变量。例如给出
$U_V|_T=F(T,V)$，积分得到
\[
 U(T,V)=\int_{V_\ast}^{V}F(T,w)\dd w+\Phi(T).
\]
再对温度求导，
\[
 C_V(T,V)=\int_{V_\ast}^{V}
 \left(\frac{\partial F}{\partial T}\right)_w\dd w+\Phi'(T).
\]
指定参考体积上的热容，才确定 $\Phi'(T)$；再给一个内能参考值，才确定最终的加法常数。若热容已经在所有体积上给出，它必须满足式\eqref{eq:cvcompat}，才能与同一内能函数相容。

熵的积分有相同结构。用一条先定容、后等温的路径，
\[
 S(T,V)-S(T_\ast,V_\ast)
 =\int_{T_\ast}^{T}\frac{C_V(t,V_\ast)}{t}\dd t
 +\int_{V_\ast}^{V}\left(\frac{\partial p(T,w)}{\partial T}\right)_w\dd w.
\]
第二项沿终温 $T$ 的等温线积分。对整个结果求温度导数，利用相容关系，可恢复终体积处的 $C_V(T,V)/T$。若选择先等温、后定容的另一条路径，相应的等温段应取初温，定容段也应改在终体积处积分；两段须一起调整。

\subsection{实际气体的绝热线与机械稳定区}
可逆绝热约束由熵微分直接给出，
\[
 \frac{C_V(T,V)}{T}\dd T+
 \left(\frac{\partial p}{\partial T}\right)_V\dd V=0.
\]
只有在系数具有可分离形式时，才可积分成简单幂律。对 $p=\nu RT/(V-B)-A/V^2$ 且总定容热容为常数的模型，
$(\partial p/\partial T)_V=\nu R/(V-B)$，由积分得到温度与有效自由体积的关系。吸引项虽然不直接出现在熵的体积系数中，却仍影响内能、绝热功和机械稳定性。

若路径进入 $(\partial p/\partial V)_T>0$ 的均匀态分支，微小密度扰动会进一步增大，单相均匀模型失稳。此时沿形式上的代数曲线继续积分，并不代表真实平衡过程；应改用两相共存条件及相比例。第四讲讨论的相变平台正是状态描述发生变化的地方。


\section{热容、状态方程与过程约束的综合应用}
状态函数的全微分联系任意相邻平衡态，过程条件则确定其中的一条路径。定压过程取 $\dd p=0$，可逆绝热过程取 $\dd S=0$，稳流节流过程取 $\dd H=0$。应用同一组状态关系，可以得到不同约束下的温度变化与热量交换。

\begin{example}[容器热容不可忽略的泄流冷却]
刚性容器的壁热容为常量 $C_w$，与容器内 $N$ 个自由度为 $f$ 的理想气体粒子保持同温 $T$。整体对外绝热，气体从小孔向真空缓慢泄出。初态 $N_0,T_0$。求 $N(T)$，并讨论极少余气时温度的形式极限。
\Diagram{fig:wallcapacityq}{气体与器壁共同储能，粒子由小孔逸出}{
 \draw[line width=2pt](0,0)rectangle(3.3,1.8);\fill[white](3.22,.73)rectangle(3.38,1.07);
 \node at(1.65,.9){$N,c_v,T$};\node at(1.65,2.3){器壁 $C_w,T$};
 \draw[axis](3.3,.9)--(5,.9);\node at(4.3,.2){$\varepsilon_e$};
}
\end{example}
\begin{solution}
图\ref{fig:wallcapacityq}中系统能量 $U=Nc_vT+C_wT$，$c_v=f\kb/2$。每个流出粒子携能 $(c_v+\kb/2)T$。取 $\dd N<0$，能量守恒为
\[
 (Nc_v+C_w)\dd T+c_vT\dd N
 =(c_v+\kb/2)T\dd N,
\]
即 $(Nc_v+C_w)\dd T=(\kb/2)T\dd N$。把 $N+C_w/c_v$ 作为整体，
\[
 \frac{\dd N}{N+C_w/c_v}=f\frac{\dd T}T.
\]
从初态积分，$[\ln(N'+C_w/c_v)]_{N_0}^{N}=f[\ln\vartheta]_{T_0}^{T}$，即$\ln[(N+C_w/c_v)/(N_0+C_w/c_v)]=f\ln(T/T_0)$。取指数得到
\[
 N+\frac{C_w}{c_v}=\left(N_0+\frac{C_w}{c_v}\right)
 \left(\frac T{T_0}\right)^f.
\]
其中$C_w/c_v$表示与器壁热容相等的气体粒子数。气体流失时，这一等效数不变，器壁所占的热容比例逐渐增大，温度下降也随之放缓。

\Diagram{fig:wallcapacitya}{器壁储能使极少余气的温度保持有限}{
 \draw[axis](0,0)--(4.7,0)node[right]{$N/N_0$};\draw[axis](0,0)--(0,2.9)node[above]{$T/T_0$};\draw[curve]plot[domain=0:1]({4*\x},{2.4*((\x+.3)/1.3)^(1/3)});\draw[guide](0,1.47)--(4,1.47);\node[right]at(4,1.47){$T_{\rm lim}/T_0$};
}
形式上 $N\to0$ 时，图\ref{fig:wallcapacitya}给 $T_{\rm lim}=T_0[C_w/(C_w+N_0c_v)]^{1/f}$。$C_w=0$ 恢复无热容器壁结果；$C_w\gg N_0c_v$ 则温降很小。末期若碰撞不足以使余气平衡，或器壁与余气的热交换太慢，这一共同温度模型会先失效，不能把形式极限当成任意稀薄程度都精确可测的终态。
\end{solution}

\begin{example}[分子气体与辐射共同膨胀]
同一腔内有 $N$ 个单原子理想气体粒子和与之平衡的辐射，温度均为 $T$。气体粒子数固定，辐射能量密度 $aT^4$，腔壁热容忽略。求总压强、总内能、熵差及可逆绝热膨胀的温度体积关系。
\Diagram{fig:gammaradq}{分子和光子共享体积与温度}{
 \draw(0,0)rectangle(4,1.8);\foreach \p in{(.6,.4),(1.2,1.2),(2.8,.5),(3.5,1.3)}\fill\p circle(1.2pt);\draw[decorate,decoration={snake,amplitude=1.2pt,segment length=8pt},->](.5,.8)--(3.4,.8);\node at(2,-.5){$N$ 固定，光子数不固定};
}
\end{example}
\begin{solution}
图\ref{fig:gammaradq}的两部分独立贡献压强和能量，
\[
 p=\frac{N\kb T}V+\frac a3T^4,\qquad
 U=\frac32N\kb T+aVT^4.
\]
总熵差是理想气体与辐射熵差之和：
\begin{align*}
 S-S_0={}&\frac32N\kb\ln\frac T{T_0}+N\kb\ln\frac V{V_0}\\
 &+\frac{4a}3(VT^3-V_0T_0^3).
\end{align*}
可逆绝热令此式为零，便得到有限变化的隐式关系。若求局部温度斜率，更直接从 $\dd U+p\dd V=0$ 出发：
\[
 \left(\frac32N\kb+4aVT^3\right)\dd T
 +\left(\frac{N\kb T}V+\frac{4a}3T^4\right)\dd V=0.
\]
熵的气体项按固定粒子数积分，辐射项按允许光子产生和吸收的平衡场积分，两者虽共享温度，粒子数约束不同。令 $X=aVT^3/(N\kb)$，图中的局部指数为 $\Gamma=(1+4X/3)/(3/2+4X)$。对 $X$ 求导，
\[
 \frac{\dd\Gamma}{\dd X}
 =\frac{\tfrac43(\tfrac32+4X)-4(1+\tfrac43X)}{(3/2+4X)^2}
 =-\frac{2}{(3/2+4X)^2}<0.
\]
因此它连续地从分子气体极限降至辐射极限。把 $X$ 当成常数仅适用于局部斜率的比较；沿实际绝热路径 $V,T$ 都变化，有限变化仍应使用完整的熵方程。

\Diagram{fig:gammarada}{局部绝热指数在两种主导机制间变化}{
 \draw[axis](0,0)--(4.9,0)node[right]{$aVT^3/(N\kb)$};\draw[axis](0,0)--(0,2.8)node[above]{$-\dd\ln T/\dd\ln V$};\draw[curve]plot[domain=0:4.4](\x,{3*(1+4*\x/3)/(1.5+4*\x)});\draw[guide](0,1)--(4.5,1);\node[left]at(0,2){$2/3$};\node[left]at(0,1){$1/3$};
}
因此图\ref{fig:gammarada}的指数为 $(N\kb+4aVT^3/3)/(3N\kb/2+4aVT^3)$。气体主导时趋于 $2/3$，辐射主导时趋于 $1/3$。中间区域不能用一个固定热容比描述整条绝热线。任意两态之间的绝热功仍由 $W=U_1-U_2$ 计算，但要先按熵约束求末温。
\end{solution}

\begin{example}[van der Waals 气体的定压热容]
一摩尔 van der Waals 气体在稳定单相区，$p=RT/(v-b)-a/v^2$，$C_{V,m}=c$。求 $C_{p,m}$，并解释绝热线的指数 $1+R/c$ 是否等于实际的 $C_{p,m}/C_{V,m}$。
\Diagram{fig:vdwcpq}{定压方向与等温线斜率的关系}{
 \draw[axis](0,0)--(4.8,0)node[right]{$v$};\draw[axis](0,0)--(0,3.2)node[above]{$p$};\draw[curve]plot[domain=1.2:4.2](\x,{2/(\x-.3)-.2/\x^2});\draw[curve]plot[domain=1.2:4.2](\x,{2.8/(\x-.3)-.2/\x^2});\draw[axis](1.6,1.460337)--(2.162827,1.460337);\node[above left]at(1.6,1.460337){$T$};\node at(3.4,2.4){$T+\dd T$};
}
\end{example}
\begin{solution}
图\ref{fig:vdwcpq}的定压方向满足 $0=p_T\dd T+p_v\dd v$，所以 $v_T|_p=-p_T/p_v$。热容差为 $C_p-C_V=Tp_Tv_T|_p=-Tp_T^2/p_v$。本题
\[
 p_T=\frac R{v-b},\qquad
 p_v=-\frac{RT}{(v-b)^2}+\frac{2a}{v^3}.
\]
故
\[
 C_{p,m}=c+\frac{TR^2/(v-b)^2}{RT/(v-b)^2-2a/v^3}
 =c+\frac R{1-2a(v-b)^2/(RTv^3)}.
\]
热容差中的分母是 $-p_v|_T$，它同时控制定压升温时所需的体积变化。分母越小，同样的温升引起的膨胀越大，定压供热与定容供热的差也越大。这里 $C_V=c$ 固定不妨碍 $C_p$ 随状态改变，因为附加项还含 $v,T$。前面绝热线中的 $1+R/c$ 由 $\dd S=c\dd T/T+R\dd v/(v-b)$ 积分得到；这一推导从未要求实际气体满足 $C_p-C_V=R$。

\Diagram{fig:vdwcpa}{稳定区趋近机械不稳定边界时热容差增大}{
 \node[draw]at(0,0){$p_v|_T<0$};\draw[axis](1.2,0)--(2.6,0);\node[draw]at(4.6,0){$C_p-C_V=-Tp_T^2/p_v>0$};\node at(2.3,-.9){$p_v\to0^-$ 时可显著增大};
}
图\ref{fig:vdwcpa}中，机械稳定要求分母为正。一般情况下$C_p-C_V$随状态变化，绝热线$(p+a/v^2)(v-b)^{1+R/c}=\text{常量}$中的指数$1+R/c$来自熵积分。当分母中的吸引修正$2a(v-b)^2/(RTv^3)$可忽略时，$C_p-C_V\simeq R$，该指数才近似等于热容比。特别地，$a=0$时这两者精确相等，即使排除体积$b$不为零；不能把这一相等关系当作理想气体的充分条件。
\end{solution}

\begin{example}[实际气体温变热容的绝热线]
一摩尔气体仍满足 van der Waals 方程，但 $C_{V,m}(T)=c_0+qT>0$。给定初态 $(v_1,T_1)$，可逆绝热膨胀至 $v_2$。写出决定 $T_2$ 的方程和做功公式，证明正温度解在适用温区内唯一。
\Diagram{fig:vdwvariableq}{绝热路径的温度斜率由即时热容决定}{
 \draw[axis](0,0)--(4.7,0)node[right]{$v$};\draw[axis](0,0)--(0,3)node[above]{$T$};\draw[curve](1,2.5)to[out=-55,in=170](4,.7);\node[above]at(1,2.5){$1$};\node[right]at(4,.7){$2$};
}
\end{example}
\begin{solution}
由于状态方程对温度仍是线性的，相容条件允许$C_V$为任意温度函数。图\ref{fig:vdwvariableq}的熵微分为$(c_0/T+q)\dd T+R\dd v/(v-b)$。取初末态作积分端点，两个原函数分别是$c_0\ln T+qT$与$R\ln(v-b)$，故
\[
 c_0\ln\frac{T_2}{T_1}+q(T_2-T_1)
 +R\ln\frac{v_2-b}{v_1-b}=0.
\]
左侧对 $T_2$ 的导数为 $c_0/T_2+q=C_V(T_2)/T_2>0$，因而在热容为正的研究温区内至多有一个解。膨胀的体积对数项为正，温度积分必须为负，所以 $T_2<T_1$。
温度积分应按两个不同的被积函数分别处理：熵中是 $(c_0+qT)/T=c_0/T+q$，内能中是 $c_0+qT$。所以前者产生对数加线性项，后者产生线性加二次项。给定 $v_2$ 后，末温方程左侧随 $T_2$ 严格增加，数值求解时只需在模型允许且热容为正的温区内寻找夹住根的两端。单调性保证至多一个解；是否存在适用范围内的解，还要由初末体积和该温区一起判断。

\Diagram{fig:vdwvariablea}{求末温是单调方程的求根问题}{
 \draw[axis](0,0)--(4.8,0)node[right]{$T_2$};\draw[axis](0,-1)--(0,2.4)node[above]{$\Delta S$};\draw[curve](.8,-.7)to[out=30,in=215](4,2);\node[below]at(2,0){根};\node at(3.4,-.6){斜率 $C_V/T_2>0$};
}
内能温度积分为$\int_{T_1}^{T_2}(c_0+qT)\dd T=[c_0T+qT^2/2]_{T_1}^{T_2}$，体积积分为$\int_{v_1}^{v_2}av^{-2}\dd v=[-a/v]_{v_1}^{v_2}$。按图\ref{fig:vdwvariablea}确定末温，绝热功取内能差的负值，得到
\[
 W=c_0(T_1-T_2)+\frac q2(T_1^2-T_2^2)-a(1/v_1-1/v_2).
\]
热容随温度变化时，末温由整个温区的熵积分决定；只取末态热容代入常热容公式会遗漏这部分变化。求出末温后，再由内能差计算功，两次积分各有明确的物理含义。
\end{solution}

\begin{example}[节流反转曲线的完整参数形式]
对一摩尔 van der Waals 气体，利用第二讲节流系数的精确表达，求反转曲线的 $T_i(v)$ 与 $p_i(v)$。讨论正压反转曲线的体积范围和最大压强。
\Diagram{fig:inversionq}{节流温变的符号由反转线分隔}{
 \draw[axis](0,0)--(5,0)node[right]{$T$};\draw[axis](0,0)--(0,3)node[above]{$p$};\draw[curve](.8,0)to[out=80,in=160](2.8,2.4)to[out=-20,in=110](4.4,0);\node at(2.7,1){冷却区示意};
}
\end{example}
\begin{solution}
图\ref{fig:inversionq}的边界满足 $Tv_T|_p=v$。由 $v_T|_p=[R/(v-b)]/[RT/(v-b)^2-2a/v^3]$，交叉相乘得
\[
 \frac{RT}{v-b}=\frac{vRT}{(v-b)^2}-\frac{2a}{v^2}
 \quad\Rightarrow\quad
 T_i(v)=\frac{2a(v-b)^2}{Rb v^2}.
\]
代入状态方程，
\[
 p_i(v)=\frac{2a(v-b)}{bv^2}-\frac a{v^2}
 =\frac{a(2v-3b)}{bv^2}.
\]
对于给定正压，体积参数可由二次方程 $p_ibv^2-2av+3ab=0$ 解出：
\[
 v_\pm=\frac{a\pm\sqrt{a^2-3ab^2p_i}}{bp_i}.
\]
因此低于最大反转压强时通常存在两个参数分支，在最大压强处合并。每个体积再代入 $T_i(v)$ 才得到相应反转温度，不能把单一稀薄反转温度外推到整条曲线。两分支的具体工况还须结合相态判断；机械局部稳定与避开两相共存区是不同的条件。

\Diagram{fig:inversiona}{反转压强作为摩尔体积的函数}{
 \draw[axis](0,0)--(5.3,0)node[right]{$v/b$};\draw[axis](0,0)--(0,2.7)node[above]{$p_i b^2/a$};\draw[curve]plot[domain=1.5:6,samples=70]({.75*\x},{6*(2*\x-3)/(\x*\x)});\draw[guide](2.25,0)--(2.25,2)--(0,2);\node[left]at(0,2){$1/3$};\node[below]at(2.25,0){$3$};
}
正压要求$v>3b/2$。将图\ref{fig:inversiona}的函数写为$p_i=(a/b)(2v^{-1}-3bv^{-2})$，逐项求导，$p_i'=(a/b)(-2v^{-2}+6bv^{-3})=(2a/b)(3b-v)/v^3$。导数在$v=3b$由正变负，故$p_{i,\max}=a/(3b^2)$，对应$T_i=8a/(9Rb)$。当$v\to\infty$，反转温度趋于$2a/(Rb)$，正是稀薄近似的结果；低体积端$v\to3b/2$则趋于$2a/(9Rb)$。

这些参数式属于 van der Waals 单相模型的数学反转线。实际运行还要排除两相区及不稳定状态，图中的几何边界本身不保证每一点都可作为稳定节流工况。
\end{solution}

\begin{example}[由低压体积展开求焓与热容]
一摩尔气体在低压区满足 $v(T,p)=RT/p+B(T)$，忽略更高阶压强项。已知零压极限摩尔焓为 $h_0(T)$。求 $h(T,p)$ 和 $C_{p,m}(T,p)$，并用 $B=b-a/(RT)$ 验证稀薄 van der Waals 节流结果。
\Diagram{fig:virialq}{低压体积相对于理想项的偏移}{
 \draw[axis](0,0)--(4.7,0)node[right]{$1/p$};\draw[axis](0,0)--(0,3)node[above]{$v$};\draw[curve](.3,.8)--(4,2.6);\draw[guide](0,.2)--(4,2.15);\node at(2,2.9){竖直偏差 $B(T)$};
}
\end{example}
\begin{solution}
图\ref{fig:virialq}的定温偏移为 $B(T)$。计算 $v_T|_p=R/p+B'(T)$，焓的等温压强导数为
\[
 h_p|_T=v-Tv_T=B-TB'.
\]
它在保留的近似阶数内与压强无关，于是从零压参考积分，$h-h_0(T)=\int_0^p[B(T)-TB'(T)]\dd p'=[p'(B-TB')]_0^p=p(B-TB')$。再对温度求导，
\[
 C_{p,m}=h_0'(T)-pTB''(T).
\]
对焓修正项求温度导数时，要同时对 $T$ 与 $B'(T)$ 求导：$\dd(TB')/\dd T=B'+TB''$，它与前一个 $B'$ 抵消，才留下 $-TB''$。零压焓 $h_0(T)$ 的导数给零压定压热容，参考内能零点只给焓增加一个常数。这里的体积式是低压展开，故所得焓保留到压强的一阶项；若要继续提高热容的压强精度，需同时提高状态方程的展开阶数。

\Diagram{fig:viriala}{同一修正函数的不同导数组合对应不同性质}{
 \node[draw,align=center]at(0,0){焓修正\\$p(B-TB')$};\node[draw,align=center]at(4.5,0){热容修正\\$-pTB''$};\node at(2.2,-1){节流分子 $Tv_T-v=TB'-B$};
}
取 $B=b-a/(RT)$，有 $B'=a/(RT^2)$，故 $TB'-B=2a/(RT)-b$，与稀薄节流公式一致。图\ref{fig:viriala}列出的各量来自不同的偏导：焓的压强修正含 $B-TB'$，再沿定压方向对温度求导，才得到热容修正中的 $-TB''$。
\end{solution}

\begin{example}[再生器怎样改变循环效率]
$\nu$ 摩尔常热容理想气体经历 Stirling 循环：在 $T_h$ 等温膨胀 $V_a\to V_b$，等容冷却到 $T_c$，低温等温压缩回 $V_a$，再等容加热到 $T_h$。比较没有再生器与理想再生器时的外界吸热和效率。体积比 $r=V_b/V_a>1$。
\Diagram{fig:stirlingq}{两等温段和两等容段组成的循环}{
 \draw[axis](0,0)--(4.7,0)node[right]{$V$};\draw[axis](0,0)--(0,3.5)node[above]{$p$};\draw[curve]plot[domain=1:3.5](\x,{3/\x});\draw[curve]plot[domain=1:3.5](\x,{1.5/\x});\draw[curve](1,1.5)--(1,3);\draw[curve](3.5,.429)--(3.5,.857);\node[below]at(1,0){$V_a$};\node[below]at(3.5,0){$V_b$};\node at(3,2.7){$T_h$};\node at(2.5,.25){$T_c$};
}
\end{example}
\begin{solution}
图\ref{fig:stirlingq}两条等容段不做功。沿等温线积分$p=\nu RT/V$，高温段的功为$W_h=\nu RT_h[\ln V]_{V_a}^{V_b}=\nu RT_h\ln r$，低温段为$W_c=\nu RT_c[\ln V]_{V_b}^{V_a}=-\nu RT_c\ln r$。相加得$W=\nu R(T_h-T_c)\ln r$。没有再生器且等容加热由外部供热时，外界总吸热
\[
 Q_h^{(0)}=\nu RT_h\ln r+C_V(T_h-T_c),
 \quad \eta_0=\frac{\nu R(T_h-T_c)\ln r}{\nu RT_h\ln r+C_V(T_h-T_c)}.
\]
两条等容段的热量分别为 $-C_V(T_h-T_c)$ 与 $+C_V(T_h-T_c)$，沿整个工质循环相加为零，但分别位于放热、吸热阶段。计算从外界取得的总正热量时，不能仅因其代数和为零就把它们事先删去。理想再生器的作用是把前一段释放的热在相应温度处储存，再交还给后一段，使这两段的热量转为装置内部交换。循环功不受再生器改变，效率差别来自外界吸热量的改变。

\Diagram{fig:stirlinga}{理想再生器在两个等容段之间暂存并返还热量}{
 \node[draw]at(0,1.3){等容冷却};\node[draw]at(0,-1.3){等容加热};\node[draw]at(4,0){再生器};\draw[axis](1,1.3)--(3.2,.3)node[midway,above]{$C_V(T_h-T_c)$};\draw[axis](3.2,-.3)--(1,-1.3);
}
图\ref{fig:stirlinga}的理想再生器沿程具有连续温度分布，各微元以无穷小温差存取热量，循环后能量和熵都恢复原值。两等容段于是改为装置内部换热，高温热源只供给$Q_h=\nu RT_h\ln r$，效率为$\eta=1-T_c/T_h$。一块均温金属通常达不到这种再生条件，因为它与各温度的气体接触时仍有有限温差。
\end{solution}

\begin{example}[零粒子流时的能量流与熵产生]
两个单原子气体大容器以分子流小孔相通，温度 $T_1>T_2$ 恒定，数密度满足零净粒子通量条件。以每单位孔面积的单向粒子通量 $j$ 表示，求净能流和两热库的总熵产生率。
\Diagram{fig:effentropyq}{粒子数双向抵消而能量不抵消}{
 \node[draw]at(0,0){$T_1,n_1$};\node[draw]at(5,0){$T_2,n_2$};\draw[axis](1.2,.4)--(3.8,.4)node[midway,above]{$j$};\draw[axis](3.8,-.4)--(1.2,-.4)node[midway,below]{$j$};
}
\end{example}
\begin{solution}
图\ref{fig:effentropyq}两边单向粒子通量相同，要求 $n_1\sqrt{T_1}=n_2\sqrt{T_2}$，并非要求两压强相同。每个从侧1逸出的粒子平均带走 $2\kb T_1$，侧2则带走 $2\kb T_2$。净能流为
\[
 J_E=2\kb j(T_1-T_2).
\]
两容器维持恒温、恒粒子数，稳态时其自身宏观熵不继续积累。热侧恒温装置需补回流出的净能量，冷侧恒温装置需取走同量能量；图中计算的正熵率属于两个维持温度的热库。由于两侧各自的净粒子流为零，状态微分 $\dd S=(\dd U-\mu\dd N)/T$ 中的净 $\mu\dd N$ 项消失。若只说“有粒子双向穿孔”就额外加入两次独立的净粒子熵项，会把流入流出抵消的部分重复计算。

\Diagram{fig:effentropya}{恒温装置从热库1向热库2转移净能量}{
 \node[draw]at(0,0){热库1熵率 $-J_E/T_1$};\node[draw]at(5,0){热库2熵率 $+J_E/T_2$};\node at(2.5,-1){净粒子交换为零};
}
按图\ref{fig:effentropya}相加两个热库的熵变化率，得到
\[
 \dot S_{\rm gen}/A_h=J_E\left(\frac1{T_2}-\frac1{T_1}\right)
 =2\kb j\frac{(T_1-T_2)^2}{T_1T_2}>0.
\]
零净粒子流使两侧粒子数保持不变，熵产生由净能流引起的热库熵变给出。若压强偏离这一条件，粒子数也会变化，还须计入粒子携带的能量及各容器化学势对应的熵变。
\end{solution}

\section{有限换热速度下的气体压缩}
\subsection{等温与绝热之间还存在相位滞后}
一份理想气体在平衡态$(p_0,V_0,T_0)$附近受小幅体积振动，壁面与恒温环境之间的总热导为$H$，总定容热容为$C_V$。忽略气体内部温差，但保留与环境的有限换热速度。由第一定律，线性温度扰动满足
\[
 C_V\dot{\delta T}+H\delta T=-p_0\dot{\delta V}.
\]
右侧使用平衡压强，因为扰动压强与体积变化率的乘积是二阶项。令$\tau=C_V/H$，体积相对扰动为$\delta V/V_0=\Re[\eta\ee^{\mathrm i\omega t}]$。齐次温度方程给$\int\dd(\delta T)/\delta T=-\int\dd t/\tau$，故自由解为$C\ee^{-t/\tau}$。对周期驱动取特解$\delta T=\Re[\widehat T\ee^{\mathrm i\omega t}]$，代入并比较复振幅，
\[
 (H+\mathrm i\omega C_V)\widehat T=-\mathrm i\omega p_0V_0\eta.
\]
总解是该特解加齐次解，常数$C$由初温减去特解在$t=0$的值确定。稳定周期中齐次项已衰减。理想气体有$p_0V_0=(\gamma-1)C_VT_0$，其中$\gamma=C_p/C_V$；以下$\delta T,\delta p$表示周期扰动的复振幅，因此
\[
 \frac{\delta T}{T_0}=-\frac{(\gamma-1)\mathrm i\omega\tau}{1+\mathrm i\omega\tau}\eta,
 \qquad \frac{\delta p}{p_0}=\frac{\delta T}{T_0}-\eta.
\]
把压强写为$\delta p=-K(\omega)\eta$，得到动态体积模量
\[
 K(\omega)=p_0\frac{1+\mathrm i\gamma\omega\tau}{1+\mathrm i\omega\tau}.
\]
以$x=\omega\tau$表示频率参数，分子分母同乘$1-\mathrm ix$，有$K=p_0[1+\gamma x^2+\mathrm i(\gamma-1)x]/(1+x^2)$，故其实部和虚部分别是
\[
 K'=p_0\left[1+\frac{(\gamma-1)x^2}{1+x^2}\right],\quad
 K''=p_0\frac{(\gamma-1)x}{1+x^2}.
\]
低频极限$K'\to p_0$对应等温压缩，高频极限$K'\to\gamma p_0$对应绝热压缩。中间频段还存在由$K''$描述的相位滞后，因而需要复响应系数；只改变实热容的数值会漏掉耗散。

\begin{example}[气体的压强—体积滞回与耗散]
按本节模型，体积作$V(t)=V_0[1+\eta_0\cos\omega t]$，$\eta_0\ll1$。求稳定周期中外界对气体所做的净功，以及频率参数使耗散最大的位置。说明气体每周期内能改变量为何为零。
\ExpansionDiagram{fig:dep-th2-volume-q}{气体由活塞周期压缩，并通过有限热导向恒温环境换热}{
\draw(0,2)--(0,0)--(3,0);\draw(0,2)--(3,2);\draw[very thick](3,0)--(3,2);\draw[<->](2.5,2.4)--(3.5,2.4);\node at(1.5,1){气体};\draw[->](1.5,0)--(1.5,-1)node[below]{$H,T_0$};\node[right]at(4.5,1){$V_0+\delta V$};}
\end{example}
\begin{solution}
图\ref{fig:dep-th2-volume-q}的稳态温度已经排除起始自由弛豫。压强实扰动为
\[
 \delta p=-\eta_0K'\cos\omega t+\eta_0K''\sin\omega t.
\]
外界净功是$W_{\rm on}=-\oint p\dd V$；常数压强项闭合积分为零，所以
\begin{align*}
 W_{\rm on}&=\int_0^{2\pi/\omega}
 \delta p\,V_0\eta_0\omega\sin\omega t\dd t\\
 &=\pi V_0\eta_0^2K''
 =\pi p_0V_0\eta_0^2(\gamma-1)\frac{x}{1+x^2}.
\end{align*}
这里令$\phi=\omega t$，$\dd\phi=\omega\dd t$，上下限为$0,2\pi$。交叉项的积分$\int_0^{2\pi}\sin\phi\cos\phi\dd\phi=[\sin^2\phi/2]_0^{2\pi}=0$；耗散项的积分$\int_0^{2\pi}\sin^2\phi\dd\phi=[\phi/2-\sin(2\phi)/4]_0^{2\pi}=\pi$，分别给出上述两项的系数。
\ExpansionDiagram{fig:dep-th2-volume-a}{压强—体积椭圆的面积对应周期耗散}{
\draw[->](0,0)--(5,0)node[right]{$V$};\draw[->](0,0)--(0,3.5)node[above]{$p$};\draw[thick,domain=0:360,samples=120]plot({2.5+1.5*cos(\x)},{1.8-.8*cos(\x)+.6*sin(\x)});\draw[->,domain=20:40,samples=12]plot({2.5+1.5*cos(\x)},{1.8-.8*cos(\x)+.6*sin(\x)});\node[below]at(2.5,0){$V_0$};\node[right]at(4,2.8){逆时针};}
图\ref{fig:dep-th2-volume-a}按气体坐标画逆时针环路，故$-\oint p\dd V>0$。用商法则，$\dd[x/(1+x^2)]/\dd x=[(1+x^2)-2x^2]/(1+x^2)^2=(1-x^2)/(1+x^2)^2$，导数在$x=1$由正变负，所以每周期耗散在该点最大。这里固定的是体积振幅，若固定驱动力振幅，响应振幅会随频率变化，极值条件需另求。

周期末温度回到同值，理想气体内能也回到同值，$\Delta U=0$。由第一定律，输入功全部变为向环境排出的净热量。环境在恒温$T_0$下的熵增加为$W_{\rm on}/T_0$，与正的耗散相符。低频时换热足够快，高频时几乎来不及换热，两端的每周期耗散均趋零。
\end{solution}

\begin{example}[热弛豫造成的活塞阻尼]
气体推动质量$M$、面积$A$的活塞，活塞还受劲度$k$的弹簧约束，外部常压在平衡点抵消$p_0$。活塞向外位移$x$使$\delta V=Ax$。在气体热修正相对总弹性很小的条件下，求相对等温固有频率的一阶频移和振幅衰减率。
\ExpansionDiagram{fig:dep-th2-piston-q}{弹簧与气体共同提供恢复力，气体还引入热耗散}{
\draw(0,0)rectangle(3,2);\draw[very thick](3,0)--(3,2);\draw(3,1)--(3.2,1)--(3.35,1.2)--(3.55,.8)--(3.75,1.2)--(3.95,.8)--(4.15,1.2)--(4.35,.8)--(4.55,1.2)--(4.75,1)--(5,1);\node[above]at(4,1.3){$k$};\draw[thick](5,0)--(5,2);\node at(1.4,1){$V_0,T_0$};\node[above]at(3,2){$M$};\draw[->](3,-.5)--(4,-.5)node[right]{$x$};}
\end{example}
\begin{solution}
图\ref{fig:dep-th2-piston-q}中压强扰动产生力$A\delta p=-A^2K(\omega)x/V_0$，所以频率方程为
\[
 -M\omega^2+k+\frac{A^2}{V_0}K(\omega)=0.
\]
等温基准频率$\omega_0^2=(k+A^2p_0/V_0)/M$。令$\omega=\omega_0+\delta\omega$、$\delta\omega=\Delta\omega+\mathrm i\Gamma$，由于采用$\ee^{\mathrm i\omega t}$，正$\Gamma$表示衰减。将频率平方展开为$\omega_0^2+2\omega_0\delta\omega$，略去二阶小量；在一阶热修正$K-p_0$内部取$\omega=\omega_0$，得到
\[
 -2M\omega_0\delta\omega+
 \frac{A^2}{V_0}[K(\omega_0)-p_0]=0.
\]
等温零阶项已由$\omega_0$的定义抵消。比较剩余式的实部与虚部，
\[
 \Delta\omega=\frac{A^2[K'(\omega_0)-p_0]}{2MV_0\omega_0},
 \qquad\Gamma=\frac{A^2K''(\omega_0)}{2MV_0\omega_0}.
\]
\ExpansionDiagram{fig:dep-th2-piston-a}{振幅随热耗散衰减，振动频率则由实部弹性决定}{
\draw[->](0,0)--(6,0)node[right]{$t$};\draw[->](0,-1.5)--(0,1.8)node[above]{$x$};\draw[dashed,domain=0:5.5]plot(\x,{1.3*exp(-.2*\x)});\draw[dashed,domain=0:5.5]plot(\x,{-1.3*exp(-.2*\x)});\draw[thick,domain=0:5.5,samples=150]plot(\x,{1.3*exp(-.2*\x)*cos(360*\x)});}
图\ref{fig:dep-th2-piston-a}示意$\ee^{-\Gamma t}$包络。令$x_0=\omega_0\tau$，显式结果为
\[
 \Delta\omega=\frac{A^2p_0(\gamma-1)}{2MV_0\omega_0}
 \frac{x_0^2}{1+x_0^2},\qquad
 \Gamma=\frac{A^2p_0(\gamma-1)}{2MV_0\omega_0}
 \frac{x_0}{1+x_0^2}.
\]
使用这些一阶式需要$A^2p_0(\gamma-1)/(MV_0\omega_0^2)\ll1$；若不满足，应回到热方程与活塞方程联立求完整特征根。仅把压缩指数选为一个实数，会遗漏$\Gamma$，也无法解释振幅逐渐减小。
\end{solution}

\section{习题与解答}
\begin{exercise}[恒温泄流所需热量]
刚性容器内有自由度为 $f$ 的理想气体，经小孔向真空泄流；外界供热使温度保持 $T_0$。求从 $N_0$ 减至 $N_1$ 所需的总热量及 $N(t)$。
\Diagram{fig:isothermalleakq}{恒温装置补热，气体经另一侧小孔泄流}{
 \draw(0,0)rectangle(3,1.7);\fill[white](2.94,.64)rectangle(3.06,.96);
 \draw[axis](-1,.8)--(0,.8)node[midway,above]{$Q$};\draw[axis](3,.8)--(4.6,.8)node[midway,above]{粒子};\node at(1.5,.8){$T_0$ 恒定};
}
\end{exercise}
\begin{solution}
图\ref{fig:isothermalleakq}中余气 $U=fN\kb T_0/2$。开放系统能量式 $\dd U=\delta Q+\varepsilon_e\dd N$ 给出
\[
 \delta Q=\frac f2\kb T_0\dd N-\frac{f+1}2\kb T_0\dd N
 =-\frac12\kb T_0\dd N.
\]
有限变化的热量由 $\int_{N_0}^{N_1}(-\kb T_0/2)\dd N$ 得到，积分下限大于上限时，负号使最终供热为正。每减少一个粒子，余气内能按体内平均能量下降，而流出粒子所带能量更高，外界只补足这个差额。即使有内部转动自由度，它们在“体内减少”和“粒子带走”两项中以相同平均值出现，所以相减后只留下半个平动热能尺度。

\Diagram{fig:isothermalleaka}{每逸出一个粒子所需补热}{
 \node[draw]at(0,0){余气减少 $f\kb T_0/2$};\node[draw]at(4.6,0){逸出带走 $(f+1)\kb T_0/2$};\node at(2.3,-.9){差额由外界供给 $\kb T_0/2$};
}
按图\ref{fig:isothermalleaka}，总热量$Q=(N_0-N_1)\kb T_0/2$，与$f$无关。通量给$\dot N=-\kappa N$，其中$\kappa=(A_h/V)\sqrt{\kb T_0/(2\pi m)}$为常量。分离变量，$\int_{N_0}^{N(t)}\dd N'/N'=-\kappa\int_0^t\dd s$，得到$\ln[N(t)/N_0]=-\kappa t$，故$N=N_0\ee^{-\kappa t}$。保持恒温仍需补热；内能因粒子减少而下降，但不足以支付全部逸出能量。
\end{solution}

\begin{exercise}[辐射腔冷却的两种热容]
体积 $V$ 的腔经小孔 $A_h$ 辐射至冷真空。分别讨论：仅计辐射场能量 $U=aVT^4$；以及腔壁热容为常量 $C_w$ 且主导总能量、腔内辐射始终与壁平衡。求两种近似的 $T(t)$。
\Diagram{fig:radcoolq}{腔壁与辐射场同温，经孔口向外辐射}{
 \draw[line width=2pt](0,0)rectangle(3.4,1.7);\fill[white](3.32,.69)rectangle(3.48,1.01);
 \node at(1.7,.85){$U_r=aVT^4$};\node at(1.7,2.2){壁热容 $C_w$};\draw[axis](3.4,.85)--(5,.85);
}
\end{exercise}
\begin{solution}
净外流功率为$A_hcaT^4/4$。若图\ref{fig:radcoolq}中只计辐射能量，能量方程为
\[
 4aVT^3\dot T=-A_hcaT^4/4.
\]
约去$aT^3$，令$\beta=A_hc/(16V)$，得到$\dot T=-\beta T$。分离变量，$\int_{T_0}^T\dd\vartheta/\vartheta=-\beta\int_0^t\dd s$，故$\ln(T/T_0)=-\beta t$，即
\[
 T=T_0\ee^{-A_hct/(16V)}.
\]

若两种热容都不能忽略，应取总热容$C_w+4aVT^3$。将能量方程除以$T^4$，并从初温积分到$T$，
\[
 \int_{T_0}^{T}\left(C_w\vartheta^{-4}+\frac{4aV}{\vartheta}\right)\dd\vartheta
 =-\frac{A_hca}{4}\int_0^t\dd s.
\]
左边的原函数为$-C_w/(3\vartheta^3)+4aV\ln\vartheta$，代入两端并同时变号，得到
\[
 \frac{C_w}{3}(T^{-3}-T_0^{-3})
 +4aV\ln\frac{T_0}{T}=\frac{A_hca}{4}t.
\]
第一项来自壁热容乘 $T^{-4}$ 的积分，第二项来自辐射热容乘 $T^{-4}$ 后剩下的 $1/T$。温度低于初温时两项都为正。分别忽略其中一个热容贡献，就恢复题设两种极限的显式温度规律。

\Diagram{fig:radcoola}{由能量温度关系决定不同的冷却微分方程}{
 \node[draw,align=center]at(0,0){辐射热容 $4aVT^3$\\$\dot T\propto-T$};
 \node[draw,align=center]at(4.5,0){腔壁热容 $C_w$\\$\dot T\propto-T^4$};
}
若壁主导，忽略混合积分中的辐射热容项，有$C_w(T^{-3}-T_0^{-3})/3=A_hcat/4$。移项并取$-1/3$次幂，得到
\[
 T=T_0\left(1+\frac{3A_hcaT_0^3}{4C_w}t\right)^{-1/3}.
\]
图\ref{fig:radcoola}表明同一个辐射功率律不决定唯一的温度时间律。第一种模型需要有能维持平衡而热容可忽略的壁；真实腔体通常要把壁与辐射热容相加，写 $(C_w+4aVT^3)\dot T=-A_hcaT^4/4$ 后再积分。
\end{solution}

\begin{exercise}[由响应系数求热容差]
某稳定单相物质的体积为 $V$，等温压缩率 $\kappa_T$、热膨胀系数 $\alpha_p$ 已知。推导 $C_p-C_V$，并对理想气体验证。
\Diagram{fig:responsq}{等压升温与等容升温的体积差}{
 \draw[axis](0,0)--(4.7,0)node[right]{$T$};\draw[axis](0,0)--(0,2.8)node[above]{$V$};\draw[axis](1,.8)--(3.5,2);\draw[axis](1,.8)--(3.5,.8);\node[right]at(3.5,2){定压};\node[right]at(3.5,.8){定容};
}
\end{exercise}
\begin{solution}
图\ref{fig:responsq}的定压路径有 $\dd V=V\alpha_p\dd T$。代入式\eqref{eq:dsvt}，再乘 $T/\dd T$，$C_p=C_V+Tp_TV\alpha_p$。状态方程全微分在定容时为 $0=V\alpha_p\dd T-V\kappa_T\dd p$，故 $p_T|_V=\alpha_p/\kappa_T$。
热容差比较的是从同一平衡态出发、同样微小温升所需的可逆热量差。定容路径直接给 $C_V\dd T$；定压路径还改变体积，其熵的额外变化由 $S_V|_T\dd V$ 给出。再用 Maxwell 关系把该偏导改写为 $p_T|_V$，便能完全由可测的响应系数表达。稳定物质 $\kappa_T>0$，故即使 $\alpha_p$ 为负，热容差仍因其平方而非负；负热膨胀本身不会令 $C_p<C_V$。

\Diagram{fig:responsa}{响应系数把状态平面中的斜率相联系}{
 \node[draw]at(0,0){$p_T|_V=\alpha_p/\kappa_T$};\draw[->](2,0)--(3,0);\node[draw]at(5,0){$C_p-C_V=TV\alpha_p^2/\kappa_T$};
}
联立两项偏导关系，得到图\ref{fig:responsa}中的$C_p-C_V=TV\alpha_p^2/\kappa_T$。理想气体$\alpha_p=1/T$、$\kappa_T=1/p$，给出$C_p-C_V=pV/T=\nu R$。若$\alpha_p=0$，两热容相等；若接近临界区响应发散，不能仍使用远离临界点的常数热容近似。
\end{solution}

\begin{exercise}[实际气体的 Joule 自由膨胀]
一摩尔 van der Waals 气体在绝热刚性大容器中从体积 $v_1$ 自由膨胀至 $v_2$，原空区为真空。$C_{V,m}=c$ 恒定。求最终温度和熵增，假设初末态均为稳定单相气体。
\Diagram{fig:jouleq}{撤去隔板后的自由膨胀}{
 \draw(0,0)rectangle(5,1.7);\draw[dashed](2,0)--(2,1.7);\node at(1,.85){$v_1,T_1$};\node at(3.5,.85){真空};\draw[->](2.2,2)--(2.2,1.4);
}
\end{exercise}
\begin{solution}
图\ref{fig:jouleq}的整体边界不移动且绝热，$Q=W=0$，所以 $U_2=U_1$。摩尔内能 $U_m=cT-a/v+\text{常量}$，故
\[
 c(T_2-T_1)+a\left(\frac1{v_1}-\frac1{v_2}\right)=0,
 \qquad T_2=T_1-\frac ac\left(\frac1{v_1}-\frac1{v_2}\right).
\]
熵增的符号还可在状态空间中直接判断。把等内能平衡态连接起来，基本关系给
\[
 0=T\dd S-p\dd V,
 \qquad \left(\frac{\partial S}{\partial V}\right)_U=\frac pT>0.
\]
这条可用于求状态差的等内能曲线，不是实际自由膨胀的准静态路径。只要它处于同一稳定单相、正温正压范围，体积增大就使熵增加。由能量式求得的末温也须先确认在此范围内，再代入熵公式。

\Diagram{fig:joulea}{实际不可逆路径与用于求熵的可逆替代路径}{
 \draw[axis](0,0)--(4.5,0)node[right]{$v$};\draw[axis](0,0)--(0,2.7)node[above]{$T$};
 \draw[axis](1,2)--(3.4,2);\draw[axis](3.4,2)--(3.4,1);\draw[guide](1,2)--(3.4,1);\node[right]at(3.4,1){$2$};\node[left]at(1,2){$1$};
}
利用图\ref{fig:joulea}的替代路径或直接用状态函数，
\[
 \Delta S=c\ln\frac{T_2}{T_1}+R\ln\frac{v_2-b}{v_1-b}.
\]
在满足模型及稳定性的可达状态间，该值非负。温度降低是吸引势能增加的结果；$a=0$ 时温度不变，但熵仍随可用体积增大。若计算给出负温度或跨入相变区，说明单相常热容模型已被用出范围。
\end{solution}

\begin{exercise}[节流系数与反转温度]
稳定流经绝热多孔塞，忽略高度和宏观动能变化。先说明为何焓不变，再对一摩尔 van der Waals 气体求 $\mu_{\rm JT}=(\partial T/\partial p)_H$，给出稀薄极限的反转温度。
\Diagram{fig:jtq}{多孔塞两侧的流动功}{
 \draw(0,0)rectangle(6,1.3);\fill[pattern=north east lines](2.8,0)rectangle(3.2,1.3);\draw[axis](.3,.65)--(1.3,.65);\draw[axis](4.6,.65)--(5.6,.65);\node at(1.5,1.8){$p_1,v_1$};\node at(4.6,1.8){$p_2,v_2$};
}
\end{exercise}
\begin{solution}
取图\ref{fig:jtq}中随一摩尔气体穿塞的能量，入口推动功为 $p_1v_1$，出口排开气体需 $p_2v_2$。绝热稳流能量守恒给 $U_{m,1}+p_1v_1=U_{m,2}+p_2v_2$，故 $H_m$ 不变。由式\eqref{eq:dhpt}，
\[
 \mu_{\rm JT}=\frac{T(\partial v/\partial T)_p-v}{C_{p,m}}.
\]
van der Waals 方程的定压微分满足
\[
 0=\frac R{v-b}\dd T+
 \left[-\frac{RT}{(v-b)^2}+\frac{2a}{v^3}\right]\dd v.
\]
因此$v_T|_p=[R/(v-b)]/[RT/(v-b)^2-2a/v^3]$。令$D_v=RT/(v-b)^2-2a/v^3>0$表示机械稳定分母，则
\[
 Tv_T-v=\frac{RT/(v-b)-vD_v}{D_v}
 =\frac{2a/v^2-RTb/(v-b)^2}{D_v},
\]
从而精确的单相结果为$\mu_{\rm JT}=[2a/v^2-RTb/(v-b)^2]/(C_{p,m}D_v)$。
低压近似体积的来源可从原方程逐项展开。对 $b/v\ll1$，$RT/(v-b)\simeq RT/v+RTb/v^2$，故 $pv\simeq RT+(RTb-a)/v$。在修正项的分母用首项 $v\simeq RT/p$ 代入，再除以 $p$，得到所用的 $v\simeq RT/p+b-a/(RT)$。这样保留的体积修正对压强是零阶，而在 $pv$ 中对应一阶低压修正，其使用范围与稀薄假设一致。

\Diagram{fig:jta}{稀薄节流时吸引与排除体积效应的竞争}{
 \draw[axis](0,0)--(5,0)node[right]{$T$};\draw[axis](0,-1)--(0,2)node[above]{$\mu_{\rm JT}$};\draw[curve]plot[domain=.9:4.5](\x,{2/\x-.8});\node[below]at(2.5,0){$T_i$};
}
稀薄近似 $v\simeq RT/p+b-a/(RT)$，所以 $Tv_T-v\simeq2a/(RT)-b$。于是
\[
 \mu_{\rm JT}\simeq\frac{2a/(RT)-b}{C_{p,m}},\qquad T_i\simeq\frac{2a}{Rb}.
\]
图\ref{fig:jta}中 $\mu_{\rm JT}>0$ 时降压 $\dd p<0$ 才导致冷却 $\dd T<0$。理想气体的焓仅依赖温度，节流不改变温度。
\end{solution}

\begin{exercise}[温度线性状态方程的内能体积项]
某物质在研究范围内满足 $p=p_*(V)+Tq(V)$。证明其内能可写为 $U=\Phi(T)-\int p_*(V)\dd V$；指出 $C_V$ 可以怎样依赖温度和体积。
\Diagram{fig:linearstateq}{在同一体积上压强随温度线性变化}{
 \draw[axis](0,0)--(4.5,0)node[right]{$T$};\draw[axis](0,0)--(0,2.7)node[above]{$p$};\draw[curve](0,.6)--(4,2.2);\node[left]at(0,.6){$p_*$};\node at(3,.7){斜率 $q(V)$};
}
\end{exercise}
\begin{solution}
在图\ref{fig:linearstateq}对应的固定体积上，$p_T=q(V)$。故 $Tp_T-p=-p_*(V)$，沿等温体积方向积分得
\[
 U(V,T)=\Phi(T)-\int^V p_*(w)\dd w.
\]
把得到的内能对温度求导时，体积积分在定容条件下不变，因而热容只有 $\Phi'(T)$。另一方面，$p$ 对温度二阶导为零，相容条件也直接要求 $C_V$ 的体积偏导为零，两种判断一致。若 $p_*(V)$ 为吸引造成的负压强项，负号再乘体积积分后可产生随膨胀上升的势能；若 $p_*$ 为零，则内能只含温度函数，理想气体正属于这一类。

\Diagram{fig:linearstatea}{内能中温度项与体积项分离}{
 \node[draw]at(0,0){$\Phi(T)$};\node at(2,0){$+$};\node[draw]at(4,0){$-\int^V p_*(w)\dd w$};\node at(2,-.8){$C_V=\Phi'(T)$，不依赖 $V$};
}
按图\ref{fig:linearstatea}求温度偏导，$C_V$ 只能依赖 $T$。熵则满足 $\dd S=\Phi'(T)\dd T/T+q(V)\dd V$。状态方程中与温度成正比的项不直接贡献内能体积项，却通过 $q(V)$ 贡献熵的体积变化。题设若再给某个随体积变化的定容热容，便与此线性状态方程不相容。
\end{solution}

\begin{exercise}[光子气体的等温膨胀]
体积由 $V_1$ 变为 $V_2$ 的光子气体与温度 $T$ 的热源接触，过程可逆。求做功、内能变化、吸热和熵增，并比较 $Q$ 与 $W$。
\Diagram{fig:photonisoq}{恒温辐射腔的可动反射壁}{
 \draw(0,0)rectangle(3.2,1.5);\draw[line width=2pt](3.2,0)--(3.2,1.5);\draw[axis](3.2,.75)--(4.5,.75);\node at(1.6,.75){$T$ 恒定};\node at(1.6,-.5){$V_1\to V_2$};
}
\end{exercise}
\begin{solution}
图\ref{fig:photonisoq}中$p=aT^4/3$不随体积变化，所以$W=\int_{V_1}^{V_2}p\dd V=p[V]_{V_1}^{V_2}=p(V_2-V_1)$。光子总数可随腔壁发射而改变，即使温度不变，能量仍随体积增加：$\Delta U=aT^4(V_2-V_1)$。
等温不变的是单位体积辐射能 $u=aT^4$，增加的体积中还要建立同样的辐射场，因此产生非零内能增量。反射壁每移动一小段，所做功为 $(u/3)\dd V$，新增辐射能为 $u\dd V$，两者比例固定。恒温热源需同时承担这两部分；增加的光子由平衡壁吸收、发射调节，其数目不受固定粒子数约束。图中的能流分配据此适用于整个可逆等温过程。

\Diagram{fig:photonisoa}{热量同时增加辐射能与支付膨胀功}{
 \node[draw]at(0,0){$Q$};\draw[->](.5,.1)--(2,1);\draw[->](.5,-.1)--(2,-1);\node[right]at(2,1){$\Delta U=3W$};\node[right]at(2,-1){$W$};
}
由第一定律 $Q=\Delta U+W=4aT^4(V_2-V_1)/3=4W$。图\ref{fig:photonisoa}中温度保持不变，辐射能仍随体积增加；其余吸热用于膨胀功。熵增 $\Delta S=Q/T=4aT^3(V_2-V_1)/3$，与状态函数表达式一致。
\end{solution}

\begingroup
\setlength{\intextsep}{3pt}\setlength{\abovecaptionskip}{4pt}
\begin{exercise}[通过熵函数反推状态方程]
某固定粒子数系统给出的熵差为 $S-S_0=C\ln(T/T_0)+\lambda\ln[(V-B)/(V_0-B)]$，其中 $C,\lambda,B$ 为常量。求与它相容的状态方程的一般形式，并说明为何还不能唯一求出压强。
\Diagram{fig:reverseq}{熵函数的两个独立方向导数}{
 \begin{scope}[yscale=.8]
 \draw[axis](0,0)--(4.4,0)node[right]{$V$};\draw[axis](0,0)--(0,2.6)node[above]{$T$};\draw[axis](1,.8)--(3,.8)node[right]{$S_V|_T$};\draw[axis](1,.8)--(1,2.1)node[above]{$S_T|_V$};
 \end{scope}
}
\end{exercise}
\begin{solution}
图\ref{fig:reverseq}中定容方向给 $C_V=TS_T|_V=C$。Maxwell 关系在等温方向给 $p_T|_V=S_V|_T=\lambda/(V-B)$。对温度积分，
\[
 p(V,T)=\frac{\lambda T}{V-B}+\psi(V).
\]
若进一步给出参考温度 $T_*$ 上的一条压强曲线 $p_*(V)$，则待定函数由 $\psi(V)=p_*(V)-\lambda T_* /(V-B)$ 唯一确定。随后内能的体积积分也随之确定，只剩一个整体参考常数。这个步骤区分了两类信息：熵给出温度变化时的压强斜率，参考等温线给出固定温度时的压强截距。两类信息合用，才足以恢复完整状态方程。

\Diagram{fig:reversea}{熵只确定压强的温度斜率}{
 \begin{scope}[yscale=.65]
 \draw[axis](0,0)--(4,0)node[right]{$T$};\draw[axis](0,0)--(0,2.8)node[above]{$p$};\draw[curve](0,.4)--(3.4,1.9);\draw[guide](0,1)--(3.4,2.5);\node at(3,.5){相同斜率，不同截距};
 \end{scope}
}
图\ref{fig:reversea}的截距 $\psi(V)$ 仍为待定函数，内能相应为 $U=CT-\int^V\psi(w)\dd w+\text{常量}$。取 $\lambda=\nu R$、$\psi=-A/V^2$ 可还原 van der Waals 模型；其他选择给出不同的物质模型。要确定 $\psi(V)$，还需一条参考等温线的压强数据，或等价的能量信息。
\end{solution}
\endgroup
